\subsection{包络代数的中心}

在本节中，我们选取 R 的一个基 B。令 \(R_{+}\) 为关于 B 的正根集合。令 \(\rho = \frac{1}{2} \sum_{\alpha \in R_{+}} \alpha\) ，并令 \(\delta\) 为代数 \(\mathbf{S}(\mathfrak{h})\) 的自同构，它将每个 \(x \in h\) 变为 \(x - \rho(x)\) ，从而将 \(h^{*}\) 上的多项式函数 p 变为函数 \(\lambda \mapsto p(\lambda - \rho)\) 。

定理 2。令 U 为 g 的包络代数，Z 为其中心，V ⊂ U 为 h 的包络代数（视为 S(h)），U⁰ 为 V 在 U 中的交换子，\(\varphi\) 为相对于 B 从 \(U^{0}\) 到 V 的 Harish-Chandra 同态（§6，第4号）。令 \(\mathbf{S}(\mathfrak{h})^{\mathrm{W}}\) 为 \(\mathbf{S}(\mathfrak{h})\) 中在 W 的作用下不变的元素所成的集合。则 \((\delta\circ\varphi)|Z\) 是从 Z 到 \(\mathbf{S}(\mathfrak{h})^{\mathrm{W}}\) 的同构，且与 B 的选择无关。

\begin{enumerate}
\def\labelenumi{\alph{enumi})}
\tightlist
\item
  令 \(P_{++}\) 为 R 的支配权集合，\(w \in W\)，\(\lambda \in P_{++}\)，\(\mu = w\lambda\)。则 \(Z(\mu - \rho)\) 同构于 \(Z(\lambda - \rho)\) 的一个子模（§6，第3号，命题6的推论2），并且对所有 \(\varphi(u)(\lambda - \rho) = \varphi(u)(\mu - \rho)\) 都有 \(u \in Z\)（§6，第4号，命题7）。因此，\((\delta \circ \varphi)(u)\) 上的多项式函数 \((\delta \circ \varphi)(u) \circ w\) 与 \(h^{*}\) 在 \(P_{++}\) 上相同。但 \(P_{++}\) 在 Zariski 拓扑下稠密于 \(h^{*}\)：这是因为可以利用由基本权 \(h^{*}\) 组成的基将 \(k^{B}\) 识别为 \(\varpi_{\alpha}\)，再应用《代数》第 IV 章第2节第3号的命题9看出这一点。因此
\end{enumerate}

\[
(\delta \circ \varphi) (u) = (\delta \circ \varphi) (u) \circ w,
\]

这证明了 \((\delta \circ \varphi)(\mathbf{Z}) \subset \mathbf{S}(\mathfrak{h})^{\mathrm{W}}\)。

\begin{enumerate}
\def\labelenumi{\alph{enumi})}
\setcounter{enumi}{1}
\tightlist
\item
  令 \(\eta\) 为第3号定理1的推论2中定义的从 \(\mathrm{I}(\mathfrak{g})\) 到 \(\mathbf{S}(\mathfrak{h})^{\mathrm{W}}\) 的同构。考虑从 g-模 U 到 g-模 \(\mathbf{S}(\mathfrak{g})\) 的典则同构（第 I 章第2节第8号），并令 \(\theta\) 为其在 Z 上的限制。则 \(\theta(\mathrm{Z}) = \mathrm{I}(\mathfrak{g})\) 。令 z 为 U 中滤子次数 \(\leq f\) 的 Z 元素。
\end{enumerate}

\[
\begin{array}{c c c} \mathrm{Z} & \stackrel {{\theta}} {{\longrightarrow}} & \mathrm{I} (\mathfrak {g}) \\ \varphi \Big \downarrow & & \Big \downarrow^ {\eta} \\ \mathbf {S} (\mathfrak {h}) & \stackrel {{\delta}} {{\longrightarrow}} & \mathbf {S} (\mathfrak {h}). \end{array}
\]

引入 §6 第4号的记号，并置

\[
z = \sum_ {\sum q _ {i} + \sum m _ {i} + \sum p _ {i} \leq f} \lambda_ {(q _ {i}), (m _ {i}), (p _ {i})} u ((q _ {i}), (m _ {i}), (p _ {i})).
\]

令 \(v((q_i),(m_i),(p_i))\) 为在 \(X_{-\alpha_1}^{q_1}\ldots X_{-\alpha_n}^{q_n}H_1^{m_1}\ldots H_l^{m_l}X_{\alpha_1}^{p_1}\ldots X_{\alpha_n}^{p_n}\) 中计算的单项式 \(\mathbf{S}(\mathfrak{g})\) 。记 \(\mathbf{S}_d(\mathfrak{g})\) 为 \(\mathbf{S}(\mathfrak{g})\) 中次数为 \(0,1,\ldots,d\) 的齐次分量之和，则有

\[
\theta (z) \equiv \sum_ {\sum q _ {i} + \sum m _ {i} + \sum p _ {i} = f} \lambda_ {(q _ {i}), (m _ {i}), (p _ {i})} v ((q _ {i}), (m _ {i}), (p _ {i})) \pmod {\mathbf {S} _ {f - 1} (\mathfrak {g})}
\]

因此

\[
(\eta \circ \theta) (z) \equiv \sum_ {\sum m _ {i} = f} \lambda_ {(0), (m _ {i}), (0)} v ((0), (m _ {i}), (0)) \pmod {\mathbf {S} _ {f - 1} (\mathfrak {h})}
\]

从而

\[
(\eta \circ \theta) (z) \equiv \varphi (z) \pmod {\mathbf {S} _ {f - 1} (\mathfrak {h})}.\tag{3}
\]

\begin{enumerate}
\def\labelenumi{\alph{enumi})}
\setcounter{enumi}{2}
\tightlist
\item
  我们证明 \(\delta \circ \varphi : \mathbf{Z} \to \mathbf{S}(\mathfrak{h})^{\mathrm{W}}\) 是双射。U 和 \(\mathbf{S}(\mathfrak{g})\) 上的典则滤子在 Z、\(\mathrm{I}(\mathfrak{g})\) 和 \(\mathbf{S}(\mathfrak{h})^{\mathrm{W}}\) 上诱导出滤子，而 \(\theta, \eta\) 与这些滤子相容，因而 \(\mathrm{gr}(\eta\circ\theta)\) 是从向量空间 \(\mathrm{gr}(\mathbf{Z})\) 到向量空间 \(\mathrm{gr}(\mathbf{S}(\mathfrak{h})^{\mathrm{W}})\) 的同构。由 (3)，\(\mathrm{gr}(\varphi)=\mathrm{gr}(\eta\circ\theta)\)，且显然 \(\mathrm{gr}(\delta)\) 是恒等映射。因此 \(\mathrm{gr}(\delta\circ\varphi)\) 是双射，所以
\end{enumerate}

\[
\delta \circ \varphi : Z \to \mathbf {S} (\mathfrak {h}) ^ {W}
\]

是双射（《交换代数》第 III 章第2节第8号定理1的推论1和2）。d）回忆 a）中的记号。令 E 为最高权为 \(\lambda\) 的单 g-模，\(\chi\) 为其中心特征标（§6，第1号，定义2）。令 \(\varphi'\) 和 \(\delta'\) 为相对于基 \(\varphi\) 而与 \(\delta\) 和 \(w(\mathrm{B})\) 类似的同态。E 相对于 \(w(\mathrm{B})\) 的最高权为 \(w(\lambda)\)。由 §6，第4号，命题7，

\[
\varphi (u) (\lambda) = \chi (u) = \varphi^ {\prime} (u) (w \lambda)
\]

对所有 \(u\in \mathbf{Z}\) 都成立，所以由 a）得，

\[
\begin{array}{c} (\delta \circ \varphi) (u) (w \lambda + w \rho) = (\delta \circ \varphi) (u) (\lambda + \rho) = \varphi (u) (\lambda) = \varphi^ {\prime} (u) (w \lambda) \\ = (\delta^ {\prime} \circ \varphi^ {\prime}) (u) (w \lambda + w \rho). \end{array}
\]

因此，多项式函数 \((\delta \circ \varphi)(u)\) 和 \((\delta' \circ \varphi')(u)\) 在 \(w(\mathrm{P}_{++}) + w\rho\) 上相同，因而相等。

推论 1。对任意 \(\lambda \in \mathfrak{h}^*\)，令 \(\chi_{\lambda}\) 为从 \(z \mapsto (\varphi(z))(\lambda)\) 到 \(Z\) 的同态 \(k\)。

\begin{enumerate}
\def\labelenumi{(\roman{enumi})}
\item
  若 \(k\) 是代数闭域，则从 \(Z\) 到 \(k\) 的每个同态都形如 \(\chi_{\lambda}\)，其中 \(\lambda \in \mathfrak{h}^*\)。
\item
  令 \(\lambda, \mu \in \mathfrak{h}^*\)。则 \(\chi_{\lambda} = \chi_{\mu}\) 当且仅当 \(\mu + \rho \in \mathrm{W}(\lambda + \rho)\)。
\end{enumerate}

若 \(k\) 是代数闭域，则从 \(\mathbf{S}(\mathfrak{h})^{\mathrm{W}}\) 到 \(k\) 的每个同态都可延拓为从 \(\mathbf{S}(\mathfrak{h})\) 到 \(k\) 的同态（《交换代数》第 V 章第1节第9号命题22，以及第2节第1号定理1的推论4），并且从 \(\mathbf{S}(\mathfrak{h})\) 到 \(k\) 的每个同态都形如 \(f \mapsto f(\lambda)\)，其中 \(\lambda \in \mathfrak{h}^*\)（第 VII 章附录 I 命题1）。因此，若 \(\chi\) 是从 \(Z\) 到 \(k\) 的同态，则存在（定理 2）某个 \(\mu \in \mathfrak{h}^*\)，使得对所有 \(z \in Z\)，

\[
\chi (z) = ((\delta \circ \varphi) (z)) (\mu) = (\varphi (z)) (\mu - \rho)
\]

因而得到 (i)。

令 \(\lambda, \mu \in \mathfrak{h}^*\)，并假设 \(\chi_{\lambda} = \chi_{\mu}\)。则对所有 \(z \in \mathbf{Z}\)，

\[
((\delta \circ \varphi) (z)) (\lambda + \rho) = (\varphi (z)) (\lambda) = \chi_ {\lambda} (z) = \chi_ {\mu} (z) = ((\delta \circ \varphi) (z)) (\mu + \rho);
\]

换言之，由 \(\mathbf{S}(\mathfrak{h})\) 和 \(\lambda + \rho\) 定义的从 \(\mu + \rho\) 到 k 的同态在 \(\mathbf{S}(\mathfrak{h})^{\mathrm{W}}\) 上相同；因此，断言 (ii) 由《交换代数》第 V 章第2节第2号定理2的推论得到。

推论 2。令 E、E' 为有限维单 g-模，\(\chi, \chi'\) 为它们的中心特征标。若 \(\chi = \chi'\)，则 E 与 E' 同构。

令 \(\lambda, \lambda'\) 为 E、E' 的最高权。由 §6，第4号，命题 7，\(\chi_{\lambda} = \chi = \chi' = \chi_{\lambda'}\)，所以存在 \(w \in \mathrm{W}\) 使 \(\lambda' + \rho = w(\lambda + \rho)\)。由于 \(\lambda + \rho\) 和 \(\lambda' + \rho\) 属于由 B 定义的房，故 \(w = 1\)。于是 \(\lambda = \lambda'\)，推论得证。

命题 7。对有限维单 g-模的任意类 \(\gamma\)，令 \(U_{\gamma}\) 为 g-模 U（取 g 在 U 上的伴随表示）中类型为 \(\gamma\) 的同构型分支。令 \(\gamma_{0}\) 为一维平凡 g-模的类。令 \([U,U]\) 为由 U 中元素对的括号生成的 U 的向量子空间。

\begin{enumerate}
\def\labelenumi{(\roman{enumi})}
\item
  U 是各 \(\mathrm{U}_{\gamma}\) 的直和。
\item
  \(U_{\gamma_{0}} = Z\)，且 \(\sum_{\gamma \neq \gamma_{0}} U_{\gamma} = [U, U]\)。
\item
  令 \(u \mapsto u^{\natural}\) 为由分解 \(\mathrm{U} = \mathrm{Z} \oplus [\mathrm{U}, \mathrm{U}]\) 定义的 U 到 Z 的投影。若 \(u \in \mathrm{U}\) 且 \(v \in \mathrm{U}\)，则 \((uv)^{\natural} = (vu)^{\natural}\)。若 \(u \in \mathrm{U}\) 且 \(z \in \mathrm{Z}\)，则 \((uz)^{\natural} = u^{\natural}z\)。
\item
  令 \(\varphi\) 为 Harish-Chandra 同态。令 \(\lambda \in \mathrm{P}_{++}\)，并令 E 为最高权为 \(\mathfrak{g}\) 的有限维单 \(\lambda\)-模。对所有 \(u \in \mathrm{U}\)，有
\end{enumerate}

\[
\frac {1}{\dim \mathrm{E}} \mathrm{Tr} (u _ {\mathrm{E}}) = (\varphi (u ^ {\natural})) (\lambda).
\]

 \(\mathfrak{g}\)-模 U 是有限维子模的直和。这蕴含 (i)。

显然 \(U_{\gamma_{0}} = Z\)。令 \(U'\) 为 U 的一个向量子空间，它在伴随表示下定义一个类为 \(\gamma\) 的子表示。于是或者 \([g, U'] = U'\)，或者 \([g, U'] = 0\)。因此，若 \(\gamma \neq \gamma_{0}\)，则 \([g, U'] = U'\)，从而 \(\sum_{\gamma \neq \gamma_{0}} U_{\gamma} \subset [U, U]\)。另一方面，若 \(u \in U\) 且 \(x_{1}, \ldots, x_{n} \in g\)，则

\[
\begin{array}{c} [ x _ {1} \ldots x _ {n}, u ] = (x _ {1} \ldots x _ {n} u - x _ {2} \ldots x _ {n} u x _ {1}) + (x _ {2} \ldots x _ {n} u x _ {1} - x _ {3} \ldots x _ {n} u x _ {1} x _ {2}) \\ + \dots + (x _ {n} u x _ {1} \ldots x _ {n - 1} - u x _ {1} \ldots x _ {n}) \in [ \mathfrak {g}, \mathrm{U} ]. \end{array}
\]

因此 \([\mathrm{U},\mathrm{U}]\subset \left[\mathfrak{g},\sum_{\gamma}\mathrm{U}_{\gamma}\right] = \left[\mathfrak{g},\sum_{\gamma \neq \gamma_0}\mathrm{U}_\gamma \right]\subset \sum_{\gamma \neq \gamma_0}\mathrm{U}_\gamma .\) 这证明了 (ii)。在这些条件下，(iii) 由第 I 章第6节第9号引理5得到。

最后，令 E、\(\lambda\) 如 (iv) 中所述。则

\[
\begin{array}{r l} & {\mathrm{Tr} (u _ {\mathrm{E}}) = \mathrm{Tr} ((u ^ {\natural}) _ {\mathrm{E}}) \qquad \text {since} u - u ^ {\natural} \in [ \mathrm{U}, \mathrm{U} ]} \\ & {\qquad = \mathrm{Tr} (\varphi (u ^ {\natural}) (\lambda). 1) \qquad (\S 6, \text {no.} 4, \text {Prop.} 7)} \\ & {\qquad = (\dim \mathrm{E}). \varphi (u ^ {\natural}) (\lambda).} \end{array}
\]

\section{H. Weyl 公式}

本段沿用 §6 和 §7 的一般记号。

\subsection{有限维 g-模的特征标}

令 \((e^{\lambda})_{\lambda \in \mathfrak{h}^{*}}\) 为环 \(\mathbf{Z}[\mathfrak{h}^*]\) 的典则基。在空间 \(\mathbf{Z}^{\mathfrak{h}^*}\) 上赋予所有从 \(\mathfrak{h}^*\) 到 \(\mathbf{Z}\) 的映射所成空间以各因子离散拓扑的积拓扑。若 \(\varphi \in \mathbf{Z}^{\mathfrak{h}^*}\)，则族 \((\varphi (\nu)e^{\nu})_{\nu \in \mathfrak{h}^*}\) 可求和，并且

\[
\varphi = \sum_ {\nu \in \mathfrak {h} ^ {*}} \varphi (\nu) e ^ {\nu}.
\]

令 \(\mathbf{Z}\langle\mathrm{P}\rangle\) 为这样的 \(\varphi\in\mathbf{Z}^{\mathfrak{h}^{*}}\) 所成的集合：其支集包含于有限个形如 \(\nu-\mathrm{P}_{+}\) 的集合之并，其中 \(\nu\in\mathfrak{h}^{*}\)。则 \(\mathbf{Z}[\mathrm{P}]\subset\mathbf{Z}\langle\mathrm{P}\rangle\subset\mathbf{Z}^{\mathfrak{h}^{*}}\)。在 \(\mathbf{Z}\langle\mathrm{P}\rangle\) 上定义一个扩展 \(\mathbf{Z}[\mathrm{P}]\) 中环结构的环结构，即对 \(\varphi,\psi\in\mathbf{Z}\langle\mathrm{P}\rangle\) 和 \(\nu\in\mathfrak{h}^{*}\)，置

\[
(\varphi \psi) (\nu) = \sum_ {\mu \in \mathfrak {h} ^ {*}} \varphi (\mu) \psi (\nu - \mu)
\]

（族 \((\varphi(\mu)\psi(\nu - \mu))_{\mu \in \mathfrak{h}^*}\) 的支集是有限的，这是因为 \(\varphi\) 和 \(\psi\) 的支集满足上述条件）。若 \(\varphi = \sum_{\nu} x_{\nu} e^{\nu}\) 且 \(\psi = \sum_{\nu} y_{\nu} e^{\nu}\)，则 \(\varphi \psi = \sum_{\nu, \mu} x_{\nu} y_{\mu} e^{\nu + \mu}\)。

令 \(\nu \in \mathfrak{h}^*\)。\(\nu\) 关于正根的一个分拆是一个族 \((n_{\alpha})_{\alpha \in \mathrm{R}_{+}}\)，其中 \(n_{\alpha}\) 为满足 \(\geq 0\) 的整数 \(\nu = \sum_{\alpha \in \mathrm{R}_{+}} n_{\alpha} \alpha\)。我们用 \(\mathfrak{P}(\nu)\) 表示 \(\nu\) 关于正根的分拆数。有

\[
\mathfrak {P} (\nu) > 0 \Longleftrightarrow \nu \in \mathrm{Q} _ {+}.
\]

本段中，用 K 表示 \(Z\langle P\rangle\) 中的以下元素：

\[
\mathrm{K} = \sum_ {\gamma \in \mathrm{Q} _ {+}} \mathfrak {P} (\gamma) e ^ {- \gamma}.
\]

现在回忆（第 VI 章第3节第3号，命题 2）：

\[
d = \prod_ {\alpha \in \mathrm{R} _ {+}} (e ^ {\alpha / 2} - e ^ {- \alpha / 2}) = \sum_ {w \in \mathrm{W}} \varepsilon (w) e ^ {w \rho}
\]

是 \(\mathbf{Z}[\mathrm{P}]\) 中的反不变元素。

引理 1。在环 \(\mathbf{Z}\langle \mathrm{P}\rangle\) 中，有 \(\mathrm{K}\) 。\(\prod_{\alpha \in \mathrm{R}_{+}}(1 - e^{-\alpha}) = \mathrm{Ke}^{-\rho}d = 1\)。事实上，

\[
\mathrm{K} = \prod_ {\alpha \in \mathrm{R} _ {+}} \left(e ^ {0} + e ^ {- \alpha} + e ^ {- 2 \alpha} + \dots\right)
\]

因此

\[
K e ^ {- \rho} d = \prod_ {\alpha \in \mathrm{R} _ {+}} \left(1 + e ^ {- \alpha} + e ^ {- 2 \alpha} + \dots\right) \prod_ {\alpha \in \mathrm{R} _ {+}} \left(1 - e ^ {- \alpha}\right) = 1.
\]

引理 2。令 \(\lambda \in \mathfrak{h}^*\)。模 \(Z(\lambda)\)（§6，第3号）有一个属于 \(\mathbf{Z}\langle P\rangle\) 的特征标，并且有 \(d.\) ch \(Z(\lambda) = e^{\lambda + \rho}\)。

令 \(\alpha_{1},\ldots ,\alpha_{q}\) 为 \(\mathbb{R}_+\) 中互不相同的元素。\(X_{-\alpha_1}^{n_1}X_{-\alpha_2}^{n_2}\dots X_{-\alpha_q}^{n_q}\otimes 1\) 构成 \(\mathrm{Z}(\lambda)\) 的一个基（§6，命题 6 (iii)）。对 \(h\in \mathfrak{h}\)，有

\[
\begin{array}{l} h. (X _ {- \alpha_ {1}} ^ {n _ {1}} X _ {- \alpha_ {2}} ^ {n _ {2}} \ldots X _ {- \alpha_ {q}} ^ {n _ {q}} \otimes 1) \\ \qquad = [ h, X _ {- \alpha_ {1}} ^ {n _ {1}} \ldots X _ {- \alpha_ {q}} ^ {n _ {q}} ] \otimes 1 + (X _ {- \alpha_ {1}} ^ {n _ {1}} \ldots X _ {- \alpha_ {q}} ^ {n _ {q}}) \otimes h. 1 \\ \qquad = (\lambda - n _ {1} \alpha_ {1} - \dots - n _ {q} \alpha_ {q}) (h) (X _ {- \alpha_ {1}} ^ {n _ {1}} \ldots X _ {- \alpha_ {q}} ^ {n _ {q}} \otimes 1). \end{array}
\]

因此，\(\mathrm{Z}(\lambda)^{\lambda - \mu}\) 的维数为 \(\mathfrak{P}(\mu)\)。这证明了 \(\operatorname{ch} \mathrm{Z}(\lambda)\) 有定义，属于 \(\mathbf{Z}\langle \mathrm{P} \rangle\)，并且

\[
\operatorname{ch} Z (\lambda) = \sum_ {\mu} \mathfrak {P} (\mu) e ^ {\lambda - \mu} = \mathrm{Ke} ^ {\lambda}.
\]

现在只需应用引理 1。

引理 3。令 M 为一个 g-模，它有一个特征标 ch(M)，其支集包含于有限个 \(\mu-P_{+}\) 之并。令 U 为 g 的包络代数，Z 为 U 的中心，\(\lambda_{0}\in h^{*}\)，\(\chi_{\lambda_{0}}\) 为相应的从 Z 到 k 的同态（§8，定理 2 的推论 1）。假设对所有 \(\S8\)\(z\in Z\)，\(z_{M}\) 都是比值为 \(\chi_{\lambda_{0}}(z)\) 的伸缩。令 \(D_{M}\) 为所有 \(\lambda\in W(\lambda_{0}+\rho)-\rho\) 所成的集合，使得 \(\lambda+Q_{+}\) 与 Supp(ch M) 相交。则 ch(M) 是所有  的 ch \(Z(\lambda)\) 的 Z-线性组合。\(\lambda\in D_{M}\)

若 \(\mathrm{Supp(chM)}\) 为空，引理显然成立。假设 \(\mathrm{Supp(chM)} \neq \emptyset\)。令 \(\lambda\) 为此支集的极大元素，并置 \(\dim M^{\lambda} = m\)。存在一个从  到 M 的 \(\mathfrak{g}\)-同态，它将 \(\varphi\)\((Z(\lambda))^{m}\)\((Z(\lambda)^{\lambda})^{m}\) 双射地映到 \(M^{\lambda}\)（§6，第3号，命题 6 (i)）。因此，\(\S 6\)\(Z(\lambda)\) 的中心特征标是 \(\chi_{\lambda_0}\)，所以 \(\lambda \in W(\lambda_0 + \rho) - \rho\)（§8，第5号，定理 2 的推论 1）。这证明 \(\S 8\)\(D_M \neq \emptyset\)，并可对 \(CardD_M\) 作归纳。令 L、N 分别为 \(\varphi\) 的核与余核。于是有如下 \(\mathfrak{g}\)-同态正合序列：

\[
0 \to \mathrm{L} \to (\mathrm{Z} (\lambda)) ^ {m} \to \mathrm{M} \to \mathrm{N} \to 0
\]

因此

\[
\operatorname{ch} (\mathrm{M}) = - \operatorname{ch} (\mathrm{L}) + m \operatorname{ch} \mathrm{Z} (\lambda) + \operatorname{ch} (\mathrm{N})
\]

（§7，第7号，公式 (6)）。Supp(ch L) 和 Supp(ch N) 包含于有限个 \(\mu -\mathrm{P}_{+}\) 之并。对 \(z\in \mathbf{Z}\)，\(z_{\mathrm{L}}\) 和 \(z_{\mathrm{N}}\) 都是比值为 \(\chi_{\lambda_{0}}(z)\) 的伸缩。显然，\(D_{N} \subset D_{M}\)。另一方面，\((\lambda + Q_{+}) \cap \text{Supp(ch M)} = \{\lambda\}\)，而 \(\lambda \notin \text{Supp(ch N)}\)，所以 \(\lambda \notin D_{N}\)，从而

\[
\mathrm{Card} \mathrm{D} _ {\mathrm{N}} <   \mathrm{Card} \mathrm{D} _ {\mathrm{M}}.
\]

另一方面，L 是 \((Z(\lambda))^{m}\) 的子模；若 \(\lambda^{\prime} \in D_{L}\)，则 \(\lambda^{\prime} + Q_{+}\) 与 Supp(ch L) ⊂ Supp ch Z(λ) 相交，所以 \(\lambda \in \lambda^{\prime} + Q_{+}\)（§6，第1号，命题 1 (ii)）；由此 \(D_{L} \subset D_{M}\)。由于 \(L \cap (Z(\lambda)^{\lambda})^{m} = 0\)，有 \(\lambda \notin D_{L}\)，所以

\(\mathrm{CardD_L} <   \mathrm{CardD_M}\)

现在只需应用归纳假设。

定理 1（H. Weyl 特征标公式）。令 M 为有限维单 g-模，\(\lambda\) 为其最高权。则

\[
\left(\sum_ {w \in \mathrm{W}} \varepsilon (w) e ^ {w \rho}\right). \mathrm{chM} = \sum_ {w \in \mathrm{W}} \varepsilon (w) e ^ {w (\lambda + \rho)}.
\]

沿用引理 3 的记号，M 的中心特征标为 \(\chi_{\lambda}\)（§6，第4号，命题 7）。因此，由引理 2 和 3，d.ch M 是满足下式的 \(\S6\)\(e^{\mu+\rho}\) 的 Z-线性组合：

\[
\mu + \rho \in W (\lambda + \rho).
\]

另一方面，由 §7，第7号，引理 7，d.ch M 是反不变的，且其唯一的极大项为 \(e^{\lambda+\rho}\)，故定理得证。

例。取 \(\mathfrak{g} = \mathfrak{sl}(2,k)\)，\(\mathfrak{h} = kH\)。令 \(\alpha\) 为 \((\mathfrak{g},\mathfrak{h})\) 的满足 \(\alpha (H) = 2\) 的根。\(\mathfrak{g}\)-模 \(V(m)\) 的最高权为 \((m / 2)\alpha\)。因此

\[
\begin{array}{r l} & {\mathrm{ch} (\mathrm{V} (m)) = (e ^ {(m / 2) \alpha + \frac {1}{2} \alpha} - e ^ {- (m / 2) \alpha - \frac {1}{2} \alpha}) / (e ^ {\frac {1}{2} \alpha} - e ^ {- \frac {1}{2} \alpha})} \\ & {\qquad = e ^ {- (m / 2) \alpha}. (e ^ {(m + 1) \alpha} - 1) / (e ^ {\alpha} - 1)} \\ & {\qquad = e ^ {- (m / 2) \alpha} (e ^ {m \alpha} + e ^ {(m - 1) \alpha} + \dots + 1)} \\ & {\qquad = e ^ {(m / 2) \alpha} + e ^ {(m - 2) \alpha / 2} + \dots + e ^ {- (m / 2) \alpha}} \end{array}
\]

这也可由 §1，第2号，命题 2 立即得到。

\subsection{单 g-模的维数}

若 \(\mu \in \mathfrak{h}^*\)，置 \(\mathrm{J}(e^{\mu}) = \sum_{w\in \mathrm{W}}\varepsilon (w)e^{w\mu}\)，参见~第 VI 章第3节第3号。

定理 2。令 E 为有限维单 g-模，\(\lambda\) 为其最高权，\((\cdot|\cdot)\) 为 \(h^{*}\) 上 W-不变的非退化正对称双线性型。则：

\[
\dim \mathrm{E} = \prod_ {\alpha \in \mathrm{R} _ {+}} \frac {\langle \lambda + \rho , H _ {\alpha} \rangle}{\langle \rho , H _ {\alpha} \rangle} = \prod_ {\alpha \in \mathrm{R} _ {+}} \left(1 + \frac {(\lambda | \alpha)}{(\rho | \alpha)}\right).
\]

令 \(\mathrm{T}\) 为不定元。对所有 \(\nu \in \mathrm{P}\)，以 \(f_{\nu}\) 表示从 \(\mathbf{Z}[\mathrm{P}]\) 到 \(\mathbf{R}[[\mathrm{T}]]\) 的同态，它对所有 \(e^{\mu}\) 将 \(e^{(\nu |\mu)\mathrm{T}}\) 映为 \(\mu \in \mathrm{P}\)。则 \(\dim \mathrm{E}\) 是级数 \(f_{\nu}(\operatorname{ch}\mathrm{E})\) 的常数项。

对所有 \(\mu, \nu \in \mathrm{P}\)，有

\[
\begin{array}{c} f _ {\nu} (\mathrm{J} (e ^ {\mu})) = \sum_ {w \in \mathrm{W}} \varepsilon (w) e ^ {(\nu | w \mu) \mathrm{T}} \\ = \sum_ {w \in \mathrm{W}} \varepsilon (w) e ^ {(w ^ {- 1} \nu | \mu) \mathrm{T}} = f _ {\mu} (\mathrm{J} (e ^ {\nu})). \end{array}
\]

特别地，根据第 VI 章第3节第3号的公式 (3)，

\[
f _ {\rho} (\mathrm{J} (e ^ {\mu})) = f _ {\mu} (\mathrm{J} (e ^ {\rho})) = e ^ {(\mu | \rho) \mathrm{T}} \prod_ {\alpha \in \mathrm{R} _ {+}} (1 - e ^ {- (\mu | \alpha) \mathrm{T}}).
\]

因此，置 \(\mathrm{Card}(\mathrm{R}_{+})=\mathrm{N}\)，

\[
f _ {\rho} (\mathrm{J} (e ^ {\mu})) \equiv \mathrm{T} ^ {\mathrm{N}} \prod_ {\alpha \in \mathrm{R}} (\mu | \alpha) \pmod {\mathrm{T} ^ {\mathrm{N} + 1} \mathbf {R} [ [ \mathrm{T} ] ]}.
\]

等式 \(\mathrm{J}(e^{\lambda +\rho}) = \mathrm{ch}(\mathrm{E}).\mathrm{J}(e^{\rho})\)（定理 1）从而蕴含

\[
\mathrm{T} ^ {\mathrm{N}} \prod_ {\alpha \in \mathrm{R} _ {+}} (\lambda + \rho | \alpha) \equiv f _ {\rho} (\operatorname{ch} \mathrm{E}). \mathrm{T} ^ {\mathrm{N}} \prod_ {\alpha \in \mathrm{R} _ {+}} (\rho | \alpha) \pmod {\mathrm{T} ^ {\mathrm{N} + 1} \mathbf {R} [ [ \mathrm{T} ] ]}
\]

因此

\[
\dim \mathrm{E} = \left(\prod_ {\alpha \in \mathrm{R} _ {+}} (\lambda + \rho | \alpha)\right) / \left(\prod_ {\alpha \in \mathrm{R} _ {+}} (\rho | \alpha)\right) = \prod_ {\alpha \in \mathrm{R} _ {+}} \left(1 + \frac {(\lambda | \alpha)}{(\rho | \alpha)}\right).
\]

现在，若 \(\alpha \in \mathrm{R}_{+}\)，则可将 \(\alpha\) 识别为 \(\mathfrak{h}_{\mathbf{R}}\) 中与 \(H_{\alpha}\) 成比例的元素，因而

\[
(\lambda + \rho | \alpha) / (\rho | \alpha) = \langle \lambda + \rho , H _ {\alpha} \rangle / \langle \rho , H _ {\alpha} \rangle .
\]

例：1）在第1号的例子中，我们得到

\[
\dim \mathrm{V} (m) = \left(\frac {m}{2} \alpha + \frac {\alpha}{2}\right) (H _ {\alpha}) / \frac {\alpha}{2} (H _ {\alpha}) = m + 1,
\]

这正是我们在 §1 中已经知道的结果。

\begin{enumerate}
\def\labelenumi{\arabic{enumi})}
\setcounter{enumi}{1}
\tightlist
\item
  取 g 为 \(G_{2}\) 型的可分裂单李代数，并采用第 VI 章图版 IX 的记号。在 \(h_{R}^{*}\) 上取 W-不变的正对称型 \((\cdot|\cdot)\)，使得 \((\alpha_{1}|\alpha_{1})=1\)。则 \(\rho=\varpi_{1}+\varpi_{2}\)，且
\end{enumerate}

\[
\begin{array}{c} (\varpi_ {1} | \alpha_ {1}) = \frac {1}{2}, (\varpi_ {1} | \alpha_ {2}) = 0, (\varpi_ {1} | \alpha_ {2} + \alpha_ {1}) = \frac {1}{2}, \\ (\varpi_ {1} | \alpha_ {2} + 2 \alpha_ {1}) = 1, (\varpi_ {1} | \alpha_ {2} + 3 \alpha_ {1}) = \frac {3}{2}, (\varpi_ {1} | 2 \alpha_ {2} + 3 \alpha_ {1}) = \frac {3}{2}, \\ (\varpi_ {2} | \alpha_ {1}) = 0, (\varpi_ {2} | \alpha_ {2}) = \frac {3}{2}, (\varpi_ {2} | \alpha_ {2} + \alpha_ {1}) = \frac {3}{2}, \\ (\varpi_ {2} | \alpha_ {2} + 2 \alpha_ {1}) = \frac {3}{2}, (\varpi_ {2} | \alpha_ {2} + 3 \alpha_ {1}) = \frac {3}{2}, (\varpi_ {2} | 2 \alpha_ {2} + 3 \alpha_ {1}) = 3. \end{array}
\]

因此，若 \(n_1, n_2\) 是 \(\geq 0\) 的整数，则最高权为 \(n_1\varpi_1 + n_2\varpi_2\) 的单表示的维数为

\[
\begin{array}{c} \left(1 + \frac {n _ {1} / 2}{\frac {1}{2}}\right) \left(1 + \frac {3 n _ {2} / 2}{\frac {3}{2}}\right) \left(1 + \frac {n _ {1} / 2 + 3 n _ {2} / 2}{\frac {1}{2} + \frac {3}{2}}\right) \left(1 + \frac {n _ {1} + 3 n _ {2} / 2}{1 + \frac {3}{2}}\right) \\ \times \left(1 + \frac {3 n _ {1} / 2 + 3 n _ {2} / 2}{\frac {3}{2} + \frac {3}{2}}\right) \left(1 + \frac {3 n _ {1} / 2 + 3 n _ {2}}{\frac {3}{2} + 3}\right) \\ = (1 + n _ {1}) (1 + n _ {2}) \left(1 + \frac {n _ {1} + 3 n _ {2}}{4}\right) \left(1 + \frac {2 n _ {1} + 3 n _ {2}}{5}\right) \left(1 + \frac {n _ {1} + n _ {2}}{2}\right) \\ \times \left(1 + \frac {n _ {1} + 2 n _ {2}}{3}\right) \\ = \frac {(1 + n _ {1}) (1 + n _ {2}) (2 + n _ {1} + n _ {2}) (3 + n _ {1} + 2 n _ {2}) (4 + n _ {1} + 3 n _ {2}) (5 + 2 n _ {1} + 3 n _ {2})}{5 !}. \end{array}
\]

特别地，最高权为 \(\varpi_{1}\)（分别为 \(\varpi_{2}\)）的基本表示的维数为 7（分别为 14）。

\subsection{单 g-模的权的重数}

命题 1。令 \(\omega \in \mathrm{P}_{++}\)。对所有 \(\lambda \in \mathrm{P}\)，\(\lambda\) 在 \(\operatorname{E}(\omega)\) 中的重数为

\[
m _ {\lambda} = \sum_ {w \in \mathrm{W}} \varepsilon (w) \mathfrak {P} (w (\omega + \rho) - (\lambda + \rho)).
\]

由定理 1 和引理 1，

\[
\operatorname{ch} \mathrm{E} (\omega) = \mathrm{K} e ^ {- \rho} d \operatorname{ch} \mathrm{E} (\omega) = \mathrm{K} e ^ {- \rho} \sum_ {w \in \mathrm{W}} \varepsilon (w) e ^ {w (\omega + \rho)}
\]

因此

\[
\operatorname{ch} \mathrm{E} (\omega) = \sum_ {w \in \mathrm{W}, \gamma \in \mathrm{Q} _ {+}} \varepsilon (w) \mathfrak {P} (\gamma) e ^ {- \rho + w (\omega + \rho) - \gamma}
\]

且

\[
m _ {\lambda} = \sum_ {w \in \mathrm{W}, \gamma \in \mathrm{Q} _ {+}, \gamma = - \lambda - \rho + w (\omega + \rho)} \varepsilon (w) \mathfrak {P} (\gamma).
\]

推论。若 \(\lambda\) 是 \(\operatorname{E}(\omega)\) 中异于 \(\omega\) 的权，

\[
m _ {\lambda} = - \sum_ {w \in \mathrm{W}, w \neq 1} \varepsilon (w) m _ {\lambda + \rho - w \rho}.
\]

在 \(\omega = 0\) 时应用命题 1。若 \(\mu \in \mathrm{P} - \{0\}\)，则得到

\[
0 = \sum_ {w \in \mathrm{W}} \varepsilon (w) \mathfrak {P} (w \rho + \mu - \rho)
\]

因而

\[
\mathfrak {P} (\mu) = - \sum_ {w \in \mathrm{W}, w \neq 1} \varepsilon (w) \mathfrak {P} (\mu + w \rho - \rho).\tag{1}
\]

命题 1 还给出

\[
m _ {\lambda} = - \sum_ {w \in W} \varepsilon (w) \sum_ {w ^ {\prime} \in W, w ^ {\prime} \neq 1} \varepsilon (w ^ {\prime}) \mathfrak {P} (w (\omega + \rho) - (\lambda + \rho) + w ^ {\prime} \rho - \rho)
\]

因为对所有 \(w(\omega + \rho) \neq \lambda + \rho\) 都有 \(w \in W\)（§7，命题 5 (iii)）。因此，

\[
\begin{array}{l} m _ {\lambda} = - \sum_ {w ^ {\prime} \in \mathrm{W}, w ^ {\prime} \neq 1} \varepsilon (w ^ {\prime}) \sum_ {w \in \mathrm{W}} \varepsilon (w) \mathfrak {P} (w (\omega + \rho) - (\lambda + \rho - w ^ {\prime} \rho + \rho)) \\ = - \sum_ {w ^ {\prime} \in \mathrm{W}, w ^ {\prime} \neq 1} \varepsilon (w ^ {\prime}) m _ {\lambda + \rho - w ^ {\prime} \rho} \quad (\text {Prop. 1}). \end{array}
\]

\subsection{单 g-模张量积的分解}

命题 2。令 \(\lambda, \mu \in \mathrm{P}_{++}\)。在 \(\mathcal{R}(\mathfrak{g})\) 中，有

\[
[ \lambda ]. [ \mu ] = \sum_ {\nu \in \mathrm{P} _ {+ +}} m (\lambda , \mu , \nu) [ \nu ]
\]

其中

\[
m (\lambda , \mu , \nu) = \sum_ {w, w ^ {\prime} \in \mathrm{W}} \varepsilon (w w ^ {\prime}) \mathfrak {P} (w (\lambda + \rho) + w ^ {\prime} (\mu + \rho) - (\nu + 2 \rho)).
\]

令 E、F 为最高权分别为 \(\lambda, \mu\) 的有限维单 g-模。令 \(l_{\nu}\) 为 E \(E \otimes F\) F 中最高权为 \(\nu\) 的同构型分支的长度。只需证明

\[
l _ {\nu} = \sum_ {w, w ^ {\prime} \in \mathrm{W}} \varepsilon (w w ^ {\prime}) \mathfrak {P} (w (\lambda + \rho) + w ^ {\prime} (\mu + \rho) - (\nu + 2 \rho)).\tag{2}
\]

置 \(c_{1} = \operatorname{ch}(\mathrm{E}) = \sum_{\sigma \in \mathrm{P}} m_{\sigma} e^{\sigma}\)，\(c_{2} = \operatorname{ch}(\mathrm{F})\)，以及 \(d = \mathrm{J}(e^{\rho})\)，其中 \(\mathrm{J}\) 如第2号所定义。有

\[
\sum_ {\xi \in \mathrm{P} _ {+ +}} l _ {\xi} \mathrm{ch} [ \xi ] = \mathrm{ch} (\mathrm{E} \otimes \mathrm{F}) = c _ {1} c _ {2}
\]

因此，在乘以 \(d\) 并应用定理1 后，

\[
\begin{array}{c} \sum_ {\xi \in \mathrm{P} _ {+ +}} l _ {\xi} \mathrm{J} (e ^ {\xi + \rho}) = c _ {1} \mathrm{J} (e ^ {\mu + \rho}) = \left(\sum_ {\sigma \in \mathrm{P}} m _ {\sigma} e ^ {\sigma}\right) \left(\sum_ {w \in \mathrm{W}} \varepsilon (w) e ^ {w (\mu + \rho)}\right) \\ = \sum_ {\tau \in \mathrm{P}} \left(\sum_ {w \in \mathrm{W}} \varepsilon (w) m _ {\tau + \rho - w (\mu + \rho)}\right) e ^ {\tau + \rho}. \end{array}\tag{3}
\]

现在，若 \(\xi \in \mathrm{P}_{++}\)，则 \(\xi + \rho\) 属于由 B 定义的房（第 VI 章第1节第10号）；因此，对所有异于 1 的 \(w \in \mathrm{W}\)，都有 \(w(\xi + \rho) \notin \mathrm{P}_{++}\)。所以，\(e^{\nu + \rho}\) 中 \(\sum_{\xi \in \mathrm{P}_{++}} l_{\xi} \mathrm{J}(e^{\xi + \rho})\) 的系数等于 \(l_{\nu}\)。由 (3) 得

\[
l _ {\nu} = \sum_ {w \in \mathrm{W}} \varepsilon (w) m _ {\nu + \rho - w (\mu + \rho)},
\]

也就是说，由命题1，

\[
l _ {\nu} = \sum_ {w, w ^ {\prime} \in \mathrm{W}} \varepsilon (w) \varepsilon (w ^ {\prime}) \mathfrak {P} (w ^ {\prime} (\lambda + \rho) - (\nu + \rho - w (\mu + \rho) + \rho))
\]

这就证明了 (2)。

例。回到第1号的例子。令 \(\lambda = (n/2)\alpha, \mu = (p/2)\alpha\)，\(\nu = (q/2)\alpha\)，其中 \(n \geq p\)。有

\[
\begin{array}{l} m (\lambda , \mu , \nu) = \mathfrak {P} \left(\frac {n}{2} \alpha + \frac {\alpha}{2} + \frac {p}{2} \alpha + \frac {\alpha}{2} - \frac {q}{2} \alpha - \alpha\right) \\ \qquad - \mathfrak {P} \left(\frac {n}{2} \alpha + \frac {\alpha}{2} - \frac {p}{2} \alpha - \frac {\alpha}{2} - \frac {q}{2} \alpha - \alpha\right) \\ \qquad - \mathfrak {P} \left(- \frac {n}{2} \alpha - \frac {\alpha}{2} + \frac {p}{2} \alpha + \frac {\alpha}{2} - \frac {q}{2} \alpha - \alpha\right) \\ \qquad + \mathfrak {P} \left(- \frac {n}{2} \alpha - \frac {\alpha}{2} - \frac {p}{2} \alpha - \frac {\alpha}{2} - \frac {q}{2} \alpha - \alpha\right) \\ \qquad = \mathfrak {P} \left(\frac {n + p - q}{2} \alpha\right) - \mathfrak {P} \left(\frac {n - p - q - 2}{2} \alpha\right). \end{array}
\]

当 \(n + p + q\) 不被 2 整除，或 \(q \geq n + p\) 时，该数为零。若

\[
q = n + p - 2 r
\]

其中 \(r\) 为 \(\geq 0\) 的整数，则有

\[
m (\lambda , \mu , \nu) = \mathfrak {P} (r \alpha) - \mathfrak {P} ((r - p - 1) \alpha)
\]

因而，当 \(m(\lambda,\mu,\nu)=1\) 时 \(r\leq p\)，当 \(m(\lambda,\mu,\nu)=0\) 时 \textgreater。最后，g-模 \(\mathrm{V}(n)\otimes\mathrm{V}(p)\) 同构于

\[
\mathrm{V} (n + p) \oplus \mathrm{V} (n + p - 2) \oplus \mathrm{V} (n + p - 4) \oplus \dots \oplus \mathrm{V} (n - p)
\]

（Clebsch–Gordan 公式）。

\section{半单李代数的极大子代数}

定理 1。令 V 为有限维向量空间，g 为 \(\mathfrak{gl}(V)\) 中的约化李子代数，q 为 g 的李子代数，\(\Phi\) 为 \((x,y)\mapsto\mathrm{Tr}(xy)\) 上的双线性型 \(g\times g\)。假设 q 关于 \(\Phi\) 的正交补 n 是 g 的李子代数，且由 V 的幂零自同态组成。则 q 是 g 的抛物子代数。

\begin{enumerate}
\def\labelenumi{\alph{enumi})}
\tightlist
\item
  \(\mathfrak{q}\) 是 \(\mathfrak{n}\) 在 \(\mathfrak{g}\) 中的正规化子：令 \(\mathfrak{p}\) 为这个正规化子。令 \(x \in \mathfrak{q}\) 且 \(y \in \mathfrak{n}\)；对所有 \(z \in \mathfrak{q}\)，都有 \([z, x] \in \mathfrak{q}\)，所以
\end{enumerate}

\[
\Phi ([ x, y ], z) = \Phi (y, [ z, x ]) = 0;
\]

换言之，\([x,y] \in n\)。因此 \(q \subset p\)。由于 n 是 p 的理想且由 V 的幂零自同态组成，P 关于 \(\varPhi\) 与 n 正交（第 I 章第3号命题4 d））。由于 \(\varPhi\) 非退化 \(^{4}\)，\(p \subset q\)，故断言得证。

\begin{enumerate}
\def\labelenumi{\alph{enumi})}
\setcounter{enumi}{1}
\tightlist
\item
  存在 \(\mathfrak{gl}(V)\) 中的约化李子代数 m，使得 q 是 m 与 n 的半直积：令 \(\mathfrak{n}_{\mathrm{V}}(\mathfrak{q})\) 为 q 中由 V 的幂零自同态组成的最大理想。则 \(\mathfrak{n}_{\mathrm{V}}(\mathfrak{q})\) 包含 n，且它与 q 正交（同上）；因此 \(\mathfrak{n} = \mathfrak{n}_{\mathrm{V}}(\mathfrak{q})\)。此外，按假设，g 是 \(\mathfrak{gl}(V)\) 中的约化李代数，因而是可分解的（第 VII 章第5节第1号命题2）；由于 q 是 g 与 \(\mathfrak{gl}(V)\) 中 n 的正规化子的交，它是一个可分解李代数（同上，命题3的推论1）。因此断言由第 VII 章第5节第3号命题7 得到。
\end{enumerate}

取 \(\mathfrak{h}\) 的一个嘉当子代数 \(\mathfrak{m}\)；记 \(\mathfrak{g}_1\) 为 \(\mathfrak{h}\) 在 \(\mathfrak{g}\) 中的交换子，并置 \(\mathfrak{q}_1 = \mathfrak{q} \cap \mathfrak{g}_1, \mathfrak{n}_1 = \mathfrak{n} \cap \mathfrak{g}_1\)。

\begin{enumerate}
\def\labelenumi{\alph{enumi})}
\setcounter{enumi}{2}
\tightlist
\item
  李代数 \(\mathfrak{g}_1, \mathfrak{q}_1\) 和 \(\mathfrak{n}_1\) 满足与 \(\mathfrak{g}, \mathfrak{q}\) 和 \(\mathfrak{n}\) 相同的假设：由于 \(\mathfrak{m}\) 是 \(\mathfrak{gl}(\mathrm{V})\) 中的约化李代数，\(\mathfrak{h}\) 是交换的，并由 V 的半单自同态组成（第 VII 章第2节第4号定理2的推论3）。因此 \(\mathfrak{g}_1 = \mathfrak{g}^0(\mathfrak{h})\) 是 \(\mathfrak{g}\) 中的约化李代数（第 VII 章第1节第3号命题11），从而也是 \(\mathfrak{gl}(\mathrm{V})\) 中的约化李代数（第 I 章第6节第6号命题7的推论2）。显然，\(\mathfrak{n}_1\) 由 V 的幂零自同态组成。由于 \(\mathfrak{h}\) 是 \(\mathfrak{q}\) 的子代数，且在 \(\mathfrak{gl}(V)\) 中约化，\(\mathfrak{h}\) 在 \(\mathfrak{q}\) 上的伴随表示是半单的；按构造，\(\mathfrak{q}_1\) 是 \(\mathrm{ad}_{\mathfrak{q}}(\mathfrak{h})\) 的不变量集合，所以 \(\mathfrak{q} = \mathfrak{q}_1 + [\mathfrak{h}, \mathfrak{q}]\)（第 I 章第3节第5号命题6）。由于
\end{enumerate}

\[
\Phi (\mathfrak {g} _ {1}, [ \mathfrak {h}, \mathfrak {q} ]) = \Phi ([ \mathfrak {h}, \mathfrak {g} _ {1} ], \mathfrak {q}) = 0,
\]

\(\mathfrak{g}_1\) 中的元素与 \(\mathfrak{q}_1\) 正交，当且仅当它与 \(\mathfrak{q}\) 正交；因此，\(\mathfrak{n}_1 = \mathfrak{g}_1 \cap \mathfrak{n}\) 是 \(\mathfrak{q}_1\) 中 \(\mathfrak{g}_1\) 的正交补。

\begin{enumerate}
\def\labelenumi{\alph{enumi})}
\setcounter{enumi}{3}
\item
   的嘉当子代数 \(\mathfrak{h}\) 是 \(\mathfrak{m}\)\(\mathfrak{g}\) 的嘉当子代数：我们有 \(\mathfrak{q} = \mathfrak{m} \oplus \mathfrak{n}\) 且 \(\mathfrak{h} = \mathfrak{m} \cap \mathfrak{g}_1\)，所以立即得到 \(\mathfrak{q}_1 = \mathfrak{h} \oplus \mathfrak{n}_1\)。此外，\([\mathfrak{h}, \mathfrak{n}_1] = 0\)，\(\mathfrak{h}\) 交换且 \(\mathfrak{n}_1\) 幂零，所以李代数 \(\mathfrak{q}_1\) 幂零。由 a）和 c），\(a)\)\(c)\)\(\mathfrak{q}_1\) 是 \(\mathfrak{n}_1\) 在 \(\mathfrak{g}_1\) 中的正规化子；更不用说，\(\mathfrak{q}_1\) 等于它在 \(\mathfrak{g}_1\) 中的正规化子，因而是 \(\mathfrak{g}_1\) 的嘉当子代数。由于 \(\mathfrak{g}_1\) 是 \(\mathfrak{gl}(V)\) 中的约化李代数，由第 VII 章第2节第4号定理2的推论3 可知，\(\mathfrak{q}_1\) 由 V 的半单自同态组成；于是，由于 \(\mathfrak{n}_1\) 由 V 的幂零自同态组成，有 \(\mathfrak{n}_1 = 0\)。因此，\(\mathfrak{h} = \mathfrak{q}_1\) 是 \(\mathfrak{g}_1\) 的嘉当子代数，并且由于 \(\mathfrak{g}_1\) 规范化 \(\mathfrak{h}\)，有 \(\mathfrak{h} = \mathfrak{g}_1\)。于是我们证明了 \(\mathfrak{h}\) 的每个元素都是 \(\mathfrak{g}\) 的半单元素，且 \(\mathfrak{h}\) 在 \(\mathfrak{g}\) 中的交换子等于 \(\mathfrak{h}\)；由此 \(\mathfrak{h} = \mathfrak{g}^0(\mathfrak{h})\)，所以 \(\mathfrak{h}\) 是 \(\mathfrak{g}\) 的嘉当子代数。
\item
  \(\mathfrak{q}\) 是 \(\mathfrak{g}\) 的抛物子代数：由上文，\(\mathfrak{h}\) 是 \(\mathfrak{g}\) 的嘉当子代数，\(\mathfrak{n}\) 由 \(\mathfrak{g}\) 的幂零元素组成，且 \([\mathfrak{h},\mathfrak{n}]\subset \mathfrak{n}\)。令 \(\bar{k}\) 为 k 的代数闭包；按定义，q 在 g 中是抛物的，当且仅当 \(k\)\(\mathfrak{q}\)\(\mathfrak{g}\)\(\bar{k}\otimes_k\mathfrak{q}\) 是 \(\bar{k}\otimes_k\mathfrak{g}\) 的抛物子代数。上述性质在标量扩张下保持不变，因此证明时可以限于 \(\mathfrak{h}\) 分裂的情形。令 R 为 \((\mathfrak{g},\mathfrak{h})\) 的根系；由第3节第1号命题2 (v)，存在 R 的子集 P，使得 \(P\cap (-P) = \varnothing\) 且 \(\mathfrak{n} = \sum_{\alpha \in P}\mathfrak{g}^{\alpha}\)。令 \(P'\) 为满足  的根 \(\alpha\) 所成的集合；有 \(-\alpha\notin P\)\(P' \cup (-P') = R\)，而  中 n 的正交补 \(\mathfrak{q}\) 等于 \(\mathfrak{n}\)\(\mathfrak{g}\)\(\mathfrak{h} + \sum_{\alpha \in P'}\mathfrak{g}^{\alpha}\)。我们已经证明 \(\mathfrak{q}\) 是抛物的。证毕。
\end{enumerate}

引理 1。令 g 为半单李代数，V 为有限维向量空间，\(\rho\) 为 g 在 V 上的线性表示，D 为 V 的向量子空间，h 为 g 的嘉当子代数，s（分别为 \(s'\)）为满足 \(x \in h\)（分别为 \(\rho(x)\mathrm{D} \subset \mathrm{D}\)）的 \(\rho(x)\mathrm{D} = 0\) 所成的集合，\(\Phi\) 为与表示 \(^{5}\) 相关联的 \(\rho\) g 上的双线性型。

\begin{enumerate}
\def\labelenumi{(\roman{enumi})}
\item
  若 \(\mathfrak{h}\) 分裂，则 \(\mathfrak{s}\) 和 \(\mathfrak{s}'\) 是 \(\mathfrak{h}\) 中关于 \(\mathbf{Q}\) 有理的向量子空间。
\item
  若 \(\rho\) 是单射，则 \(\Phi\) 在 \(\mathfrak{s}\)（分别为 \(\mathfrak{s}'\)）上的限制是非退化的。
\end{enumerate}

假设嘉当子代数 \(\mathfrak{h}\) 是分裂的。令 \(d\) 为 D 的维数；置 \(\mathrm{W} = \bigwedge^{d}(\mathrm{V})\) 和 \(\sigma = \bigwedge^{d}(\rho)\)；还以 \((e_1, \ldots, e_d)\) 表示 D 的一个基，并令 \(e = e_1 \wedge \cdots \wedge e_d\) 为与 D 相关的可分解 \(d\)-向量。令 P 为 \(\sigma\) 关于 \(\mathfrak{h}\) 的权集合；记 \(\mathrm{W}^{\mu}\) 为

与权 \(\mu\) 对应的 W 的子空间，并置 \(e = \sum_{\mu \in P} e^{\mu}\)（其中对所有 ，均有 \(e^{\mu} \in W^{\mu}\)）；最后，令 P' 为满足 \(\mu \in P\)\(P'\) 的权 \(\mu\) 所成的集合，令 P'' 为 P' 中元素的差集。令 x 属于 h；则 x 属于 s，当且仅当存在 c ∈ k 使 \(e^{\mu} \neq 0\)\(P''\)\(P'\)\(\rho(x).e = c.e\)（第 VII 章第5节第4号引理2 (i)）。由于 \(\rho(x).e^{\mu} = \mu(x).e^{\mu}\)，可见 \(x \in s\) 等价于关系“对所有 ，\(\mu(x) = 0\)”。现在，h 的 Q-结构是由余根 \(\mu \in P''\) 生成的 h 的 Q-向量子空间 \(h_{Q}\)，且 P'' 中的所有 \(H_{\alpha}\)\(\mu\) 在 \(P''\)\(h_{Q}\) 上取有理值；由此（《代数》第 II 章第8节第4号命题5），s 是 h 中关于 Q 有理的子空间。

对任意权 \(\mu\in P\)，令 \(p_{\mu}\) 为与分解  相对应的到 \(V^{\mu}\) 的投影；记 \(V=\bigoplus_{\mu\in P}V^{\mu}\)\(P_{1}\) 为满足  的 \(\mu\in P\) 所成的集合。显然，s' 是 h 中 P＿\(p_{\mu}(D)\neq0\)\(s'\)\(P_{1}\)\(s'\)1 各元素的核的交；如同 s 的情形一样可知，s' 是 h 中关于 Q 有理的子空间。这证明了 (i)。

通过标量扩张，只需在 k 为代数闭域、因而 \(k\)\(\mathfrak{h}\) 分裂时证明 (ii)。令 \(\mathfrak{m}\) 为 h 中关于 \(\mathfrak{h}\)\(\mathbf{Q}\) 有理的向量子空间；对 \(x\)\(\mathfrak{m}_{\mathbf{Q}} = \mathfrak{m} \cap \mathfrak{h}_{\mathbf{Q}}\) 中的每个非零 x，由 §7 第1号命题1的推论，有 \(\varPhi(x, x) > 0\)。\(\varPhi\) 在 \(\mathfrak{m}_{\mathbf{Q}}\) 上的限制非退化，从而由于  与  典则同构，\(\varPhi\) 在 \(\mathfrak{m}\) 上的限制也非退化。\(\mathfrak{m}\)\(k \otimes_{\mathbf{Q}} \mathfrak{m}_{\mathbf{Q}}\)

定义 1。令 q 为半单李代数 g 的一个李子代数。若对每个 \(x \in q\)，x 在 g 中的半单分量和幂零分量都属于 q，则称 q 在 g 中可分解。记 \(\mathfrak{n}_{\mathfrak{g}}(\mathfrak{q})\) 为 q 的根基中满足 \(ad_{g}x\) 幂零的元素 x 所成的集合。

令 \(\rho\) 为 g 在有限维向量空间 V 上的一个单射表示。我们知道（第 I 章第6节第3号定理3），g 的元素 x 是半单的（分别为幂零的），当且仅当 V 的自同态 \(\mathfrak{g}\)\(x\)\(\mathfrak{g}\)\(\rho(x)\) 是半单的（分别为幂零的）。由此立即得到，代数 \(\mathfrak{q}\) 在 \(\mathfrak{g}\) 中可分解，当且仅当 \(\rho(\mathfrak{q})\) 是 \(\mathfrak{gl}(V)\) 中按第 VII 章第5节第1号定义1 所称的可分解子代数。沿用第 VII 章第5节第3号的记号，还有

\[
\rho (\mathfrak {n} _ {\mathfrak {g}} (\mathfrak {q})) = \mathfrak {n} _ {\mathrm{V}} (\rho (\mathfrak {q})).
\]

定理 2。令 g 为半单李代数，n 为由幂零元素组成的 g 的子代数，q 为 n 在 g 中的正规化子。假设 n 正好是 q 的根基中的幂零元素所成的集合。则 q 是抛物的。

首先注意，q 是可分解的（第 VII 章第5节第1号命题3的推论1）。由定理1，只需证明 q 是关于 g 的 Killing 型 \(n^{0}\) 的 n 的正交补 \(\varPhi\)。我们知道 \(q \subset n^{0}\)（第 I 章第4节第3号命题4 d））。由第 VII 章第5节第3号命题7，存在 q 的一个在 g 中约化的子代数 m，使得 q 是 m 与 n 的半直积。~我们证明 \(\varPhi\) 在 m 上的限制非退化。令 c 为 m 的中心。由第 I 章第5节第5号命题5，有 \(\Phi([\mathfrak{m},\mathfrak{m}],\mathfrak{c})=0\)；而由第 I 章第6节第1号命题1，\(\Phi\) 在 \([m,m]\) 上的限制非退化。还需证明 \(\Phi\) 在 c 上的限制非退化。令 k 为 \([m,m]\) 的一个嘉当子代数；于是 \(k\oplus c\) 是交换的且在 g 中约化。令 h 为包含 \(k\oplus c\) 的 g 的一个嘉当子代数（第 VII 章第2节第3号命题10）。于是 \(h\cap q\) 是 q 的一个交换子代数，包含 \(k\oplus c\)，并且对每个 \(ad_{q}x\)，\(x\in h\cap q\) 是半单的；因此 \(h\cap q\) 包含于 q 的一个嘉当子代数 \(h'\) 中（第 VII 章第2节第3号命题10）；令 f 为以 n 为核的 q 到 m 的投影；则 \(f(h')\) 是 m 的一个嘉当子代数（第 VII 章第2节第1号命题4的推论2），包含 \(k\oplus c\)，所以等于 \(k\oplus c\)；这证明了 \(f(h\cap q)=k\oplus c\)，并且由于 h 的每个元素在 g 中都是半单的，有 \(h\cap q=k\oplus c\)。于是，

\[
\mathfrak {c} = \{x \in \mathfrak {h} \mid [ x, \mathfrak {n} ] \subset \mathfrak {n} \text {   and   } [ x, [ \mathfrak {m}, \mathfrak {m} ] ] = 0 \}.
\]

由引理1，\(\varPhi\) 在 \(\mathfrak{c}\) 上的限制非退化。

令 \(\mathfrak{q}^0\) 为 \(\mathfrak{q}\) 中相对于 \(\mathfrak{g}\) 的 \(\varPhi\) 的正交补。上文证明了 \(\mathfrak{q} \cap \mathfrak{q}^0 = \mathfrak{n}\)。假设 \(\mathfrak{q} \neq \mathfrak{q}^0\)，于是 \(\mathfrak{q}^0 \neq \mathfrak{n}\)（且 \(\mathfrak{q}^0 \supset \mathfrak{n}\)）。由于 \(\operatorname{ad}_{\mathfrak{g}} \mathfrak{n}\) 保持 \(\mathfrak{q}\) 不变，\(\operatorname{ad}_{\mathfrak{g}} \mathfrak{n}\) 也保持 \(\mathfrak{q}^0\) 不变；Engel 定理说明存在 \(x \in \mathfrak{q}^0\)，使得 \(x \notin \mathfrak{n}\) 且 \([x, \mathfrak{n}] \subset \mathfrak{n}\)。但此时 \(x \in \mathfrak{q}^0 \cap \mathfrak{q} = \mathfrak{n}\)，矛盾。因此 \(\mathfrak{q} = \mathfrak{n}^0\)。

推论 1。令 \(\mathfrak{q}\) 为异于 \(\mathfrak{g}\) 的 \(\mathfrak{g}\) 的子代数集合中的极大元素。则 \(\mathfrak{q}\) 或者是抛物的，或者在 \(\mathfrak{g}\) 中约化。

我们可以假设 g 是某个有限维向量空间 V 的 \(\mathfrak{gl}(V)\) 子代数。令 \(\mathfrak{e}(\mathfrak{q}) \subset \mathfrak{g}\) 为 q 的可分解包络。若 \(\mathfrak{e}(\mathfrak{q}) = \mathfrak{g}\)，则 q 是 g 的理想（第 VII 章第5节第2号命题4），因而是半单的，从而 q 在 g 中约化。假设 \(\mathfrak{e}(\mathfrak{q}) \neq \mathfrak{g}\)。则 \(\mathfrak{e}(\mathfrak{q}) = \mathfrak{q}\)，所以 q 可分解。假设 q 在 g 中不约化。令 n 为 q 的根基中的幂零元素所成的集合。则 \(n \neq 0\)（第 VII 章第5节第3号命题7 (i)）。令 p 为 n 在 g 中的正规化子。则 \(p \supset q\)，且由于 g 是半单的，\(p \neq g\)。于是 p = q。因此 q 是抛物的（定理1）。

推论 2。令 \(\mathfrak{n}\) 为由幂零元素组成的 \(\mathfrak{g}\) 的子代数。存在 \(\mathfrak{q}\) 的一个抛物子代数 \(\mathfrak{g}\)，具有以下性质：

\begin{enumerate}
\def\labelenumi{(\roman{enumi})}
\item
  \(\mathfrak{n}\subset \mathfrak{n}_{\mathfrak{g}}(\mathfrak{q});\)
\item
  \(\mathfrak{n}\) 在 \(\mathfrak{g}\) 中的正规化子包含于 \(\mathfrak{q}\)；
\item
  每个保持 n 不变的 g 的自同构也保持 q 不变。
\end{enumerate}

若 \(\mathfrak{g}\) 可分裂，则 \(\mathfrak{n}\) 包含于 \(\mathfrak{g}\) 的一个 Borel 子代数中。

令 \(\mathfrak{q}_1\) 为 \(\mathfrak{n}\) 在 \(\mathfrak{g}\) 中的正规化子。这是 \(\mathfrak{g}\) 的可分解子代数。令 \(\mathfrak{n}_1 = \mathfrak{n}_{\mathfrak{g}}(\mathfrak{q}_1)\)。归纳地定义 \(\mathfrak{q}_i\) 为 \(\mathfrak{n}_{i-1}\) 在 \(\mathfrak{g}\) 中的正规化子，并令 \(\mathfrak{n}_i\) 等于 \(\mathfrak{n}_{\mathfrak{g}}(\mathfrak{q}_i)\)。序列 \((\mathfrak{n}, \mathfrak{n}_1, \mathfrak{n}_2, \ldots)\) 和 \((\mathfrak{q}_1, \mathfrak{q}_2, \ldots)\) 都是递增的。存在 \(j\) 使 \(\mathfrak{q}_j = \mathfrak{q}_{j+1}\)，换言之，\(\mathfrak{q}_j\) 是 \(\mathfrak{n}_{\mathfrak{g}}(\mathfrak{q}_j)\) 在 \(\mathfrak{g}\) 中的正规化子。因此 \(\mathfrak{q}_j\) 是抛物的（定理1）。我们有 \(\mathfrak{n} \subset \mathfrak{n}_j =\) \(\mathfrak{n}_{\mathfrak{g}}(\mathfrak{q}_{j})\)，且 \(q_{1} \subset q_{j}\)；每个保持 n 不变的 g 的自同构显然都保持 \(n_{1}, n_{2}, \ldots\) 和 \(q_{1}, q_{2}, \ldots\) 不变。若 g 可分裂，则 \(q_{j}\) 包含一个 Borel 子代数 b，从而（§3，第4号，命题13），有 \(\mathfrak{b} \supset \mathfrak{n}_{\mathfrak{g}}(\mathfrak{q}_{j}) \supset \mathfrak{n}\)。

定理 3。假设 k 是代数闭域。令 g 为半单李代数，a 为 g 的可解子代数。则存在包含 a 的 g 的一个 Borel 子代数。

由第 VII 章第5节第2号命题4的推论1 (ii)，可假设 \(\mathfrak{a}\) 是可分解的。存在 g 的一个交换子代数 \(\mathfrak{t}\)，由半单元素组成，使得 \(\mathfrak{g}\)\(\mathfrak{a}\) 是 \(\mathfrak{t}\) 与 \(\mathfrak{n}_{\mathfrak{g}}(\mathfrak{a})\) 的半直积（第 VII 章第5节第3号命题6的推论2）。存在（定理2的推论2）g 的一个抛物子代数 \(\mathfrak{q}\)，使得 \(\mathfrak{g}\)\(\mathfrak{n}_{\mathfrak{g}}(\mathfrak{a}) \subset \mathfrak{n}_{\mathfrak{g}}(\mathfrak{q})\)，且 \(\mathfrak{n}_{\mathfrak{g}}(\mathfrak{a})\) 在 g 中的正规化子包含于 \(\mathfrak{g}\)\(\mathfrak{q}\)；更不用说，\(\mathfrak{a} \subset \mathfrak{q}\)。令 \(\mathfrak{b}\) 为包含于 \(\mathfrak{g}\)\(\mathfrak{q}\) 的 g 的一个 Borel 子代数，令 \(\mathfrak{h}\) 为包含于 \(\mathfrak{g}\)\(\mathfrak{b}\) 的 g 的一个嘉当子代数。于是 \(\mathfrak{h}\) 是 \(\mathfrak{q}\) 的嘉当子代数，所以存在 \(s \in \mathrm{Aut}_e(\mathfrak{q})\)，使 \(s(\mathfrak{t}) \subset \mathfrak{h}\)（第 VII 章第2节第3号命题10和第 VII 章第3节第2号定理1）。我们有 \(s(\mathfrak{n}_{\mathfrak{g}}(\mathfrak{q})) = \mathfrak{n}_{\mathfrak{g}}(\mathfrak{q})\)（第 VII 章第3节第1号注），所以

\[
s (\mathfrak {a}) = s (t) + s \left(\mathfrak {n} _ {\mathfrak {g}} (\mathfrak {a})\right) \subset \mathfrak {h} + s \left(\mathfrak {n} _ {\mathfrak {g}} (\mathfrak {q})\right) = \mathfrak {h} + \mathfrak {n} _ {\mathfrak {g}} (\mathfrak {q}) \subset \mathfrak {b}.
\]

推论。若 \(k\) 是代数闭域，则 \(\mathfrak{g}\) 的每个极大可解子代数都是 Borel 子代数。

\section{\texorpdfstring{幂零元类与 \(\mathfrak{sl}_2\) 三元组}{\textbackslash mathfrak\{sl\}\_2 三元组}}

本段中，g 表示一个有限维李代数。

\subsection{\texorpdfstring{\(\mathfrak{sl}_2\) 三元组的定义}{\textbackslash mathfrak\{sl\}\_2 三元组的定义}}

定义 1。g 中的一个 \(\mathfrak{sl}_2\) 三元组是 g 中元素的一个序列 \(\mathfrak{g}\)\((x, h, y)\)\(\mathfrak{g}\)\((0, 0, 0)\)，它异于 ，并满足

\[
[ h, x ] = 2 x, [ h, y ] = - 2 y, [ x, y ] = - h.
\]

令 \((x,h,y)\) 为 g 中的一个 \(\mathfrak{sl}_{2}\) 三元组。从 \(\tau\) 到 g 的线性映射 \(\mathfrak{sl}(2,k)\)，满足 \(\tau(X_{+}) = x, \tau(H) = h, \tau(X_{-}) = y\)，是一个非零因而为单射的同态（因为 \(\mathfrak{sl}(2,k)\) 是单的），其像为 \(kx + kh + ky\)。于是我们得到从 g 中 \(\mathfrak{sl}_{2}\) 三元组的集合到从 \(\mathfrak{sl}(2,k)\) 到 g 的单射同态集合的一个典则双射。若 g 是半单的，且 \((x,h,y)\) 是 g 中的一个 \(\mathfrak{sl}_{2}\) 三元组，则 x、y 是 g 的幂零元素，而 h 是 g 的半单元素（第 I 章第6节第3号命题4）。

引理 1。令 \(x, h, y, y' \in \mathfrak{g}\)。若 \((x, h, y)\) 和 \((x, h, y')\) 是 \(\mathfrak{sl}_2\) 中的 \(\mathfrak{g}\) 三元组，则 \(y = y'\)。

事实上，\(y - y' \in \mathrm{Ker}(\mathrm{ad}_{\mathfrak{g}}x)\)，且 \((\mathrm{ad}_{\mathfrak{g}}h)(y - y') = -2(y - y')\)。但对于每个整数 \(\mathrm{ad}_{\mathfrak{g}}x\)，\(\mathrm{Ker}(p + \mathrm{ad}_{\mathfrak{g}}h)\) 在 \(p > 0\) 上都是单射（§1，第2号，命题2的推论）。

引理 2。令 \(\mathfrak{n}\) 为 \(\mathfrak{g}\) 的一个子代数，使得对每个 \(n\in \mathfrak{n}\)，\(\operatorname{ad}_{\mathfrak{g}}(n)\) 都是幂零的。令 \(h\in \mathfrak{g}\) 满足 \([h,\mathfrak{n}] = \mathfrak{n}\)。则 \(e^{\operatorname{ad}_{\mathfrak{g}}\mathfrak{n}}.h = h + \mathfrak{n}\)。

显然 \(e^{\operatorname{ad}_{\mathfrak{g}}(\mathfrak{n})}.h \subset h + \mathfrak{n}\)。我们将证明，若 \(v \in n\)，则 \(h + v \in e^{\operatorname{ad}_{\mathfrak{g}}(\mathfrak{n})}.h\)。只需证明对所有 \(h + v \in e^{\operatorname{ad}_{\mathfrak{g}}(\mathfrak{n})}.h + \mathscr {C} ^ {p} \mathfrak {n}\)，\(p \geq 1\)（因为当 p 足够大时 \(C^{p}n = 0\)）。当 p = 1 时这显然成立，因为 \(C^{1}n = n\)。现在假设我们已经证明存在 \(y_{p} \in n\) 和 \(z_{p} \in C^{p}n\)，使得 \(h + v = e^{\operatorname{ad}_{\mathfrak{g}}y_{p}}.h + z_{p}\)。由于 \((\operatorname{ad}_{\mathfrak{g}}h)(\mathfrak{n}) = \mathfrak{n}\)，\((\operatorname{ad}_{\mathfrak{g}}h)|\mathfrak{n}\) 是从 n 到 n 的双射；因此，它在保持 \(C^{p}n\) 不变的 \(C^{p}n\) 上的限制也是双射；于是存在 \(z \in C^{p}n\)，使得 \(z_{p} = [z, h]\)。于是

\[
e ^ {\operatorname{ad} _ {\mathfrak {g}} \left(y _ {p} + z\right)} h - e ^ {\operatorname{ad} _ {\mathfrak {g}} y _ {p}} h \in [ z, h ] + \mathscr {C} ^ {p + 1} \mathfrak {n}
\]

因此

\[
e ^ {\operatorname{ad} _ {\mathfrak {g}} \left(y _ {p} + z\right)} h \in h + v - z _ {p} + [ z, h ] + \mathscr {C} ^ {p + 1} \mathfrak {n} = h + v + \mathscr {C} ^ {p + 1} \mathfrak {n}
\]

这就完成了关于 p 的归纳证明。

引理 3。令 \(x \in \mathfrak{g}\)，\(\mathfrak{p} = \operatorname{Ker}(\operatorname{ad}x)\)，\(\mathfrak{q} = \operatorname{Im}(\operatorname{ad}x)\)。则 \([\mathfrak{p}, \mathfrak{q}] \subset \mathfrak{q}\)，且 \(\mathfrak{p} \cap \mathfrak{q}\) 是 \(\mathfrak{g}\) 的子代数。

若 \(u \in \mathfrak{p}\) 且 \(v \in \mathfrak{q}\)，则存在 \(w \in \mathfrak{g}\) 使 \(v = [x, w]\)，所以

\[
[ u, v ] = [ u, [ x, w ] ] = [ x, [ u, w ] ] - [ [ x, u ], w ] = [ x, [ u, w ] ] \in \mathfrak {q}.
\]

另一方面，p 是 g 的子代数，所以 \([p \cap q, p \cap q] \subset p \cap q\)。

引理 4。令 \((x,h,y)\) 和 \((x,h',y')\) 为 \(\mathfrak{sl}_2\) 中的 \(\mathfrak{g}\) 三元组。存在 \(z\in \mathfrak{g}\)，使得 \(\mathrm{ad}_{\mathfrak{g}}z\) 幂零，并且

\[
e ^ {\mathrm{ad} _ {\mathfrak {g}} z} x = x, \quad e ^ {\mathrm{ad} _ {\mathfrak {g}} z} h = h ^ {\prime}, \quad e ^ {\mathrm{ad} _ {\mathfrak {g}} z} y = y ^ {\prime}.
\]

令 \(\mathfrak{n} = \mathrm{Ker}(\mathrm{ad} x) \cap \mathrm{Im}(\mathrm{ad} x)\)。对所有 \(p \in \mathbf{Z}\)，令 \(\mathfrak{g}_p = \mathrm{Ker}(\mathrm{ad} h - p)\)。由 §1，第3号（应用于 \(kx + ky + kh\) 在 \(\mathfrak{g}\) 上的伴随表示），有 \(\mathfrak{n} \subset \sum_{p > 0} \mathfrak{g}_p\)，所以对所有 \(\mathrm{ad}_{\mathfrak{g}} n\)，\(n \in \mathfrak{n}\) 都是幂零的，且 \([h, \mathfrak{n}] = \mathfrak{n}\)。我们有 \([x, h' - h] = 0\) 以及 \([x, y - y'] = h' - h\)，所以 \(h' - h \in \mathfrak{n}\)。由引理2和3，存在 \(z \in \mathfrak{n}\)，使 \(e^{\mathrm{ad}_{\mathfrak{g}} z} h = h'\)。由于 \(z \in \mathrm{Ker} \mathrm{ad}_{\mathfrak{g}} x\)，有 \(e^{\mathrm{ad}_{\mathfrak{g}} z} x = x\)。现在由引理1 得 \(e^{\mathrm{ad}_{\mathfrak{g}} z} y = y'\)。证毕。

令 G 为 g 的自同构群。若存在  使 、、，则称 g 中的两个 \(sl_{2}\) 三元组 \((x,h,y)\)、\((x',h',y')\) G-共轭。\(g\in G\)\(gx=x'\)\(gh=h',gy=y'\)

命题 1。令 G 为包含 \(\mathrm{Aut}_e(\mathfrak{g})\) 的 g 的自同构群。令 \((x,h,y)\) 和 \((x',h',y')\) 为 g 中的 \(\mathfrak{sl}_2\) 三元组。令

\[
\mathfrak {t} = k x + k h + k y, \quad \mathfrak {t} ^ {\prime} = k x ^ {\prime} + k h ^ {\prime} + k y ^ {\prime}.
\]

考虑以下条件：

\begin{enumerate}
\def\labelenumi{(\roman{enumi})}
\item
  x 与 \(x'\) G-共轭；
\item
  \((x, h, y)\) 与 \((x', h', y')\) G-共轭；
\item
  \(\mathfrak{t}\) 与 \(\mathfrak{t}'\) G-共轭。
\end{enumerate}

有 (i) \(\Longleftrightarrow\) (ii) \(\Longrightarrow\) (iii)。若 \(k\) 是代数闭域，则这三个条件等价。

\begin{enumerate}
\def\labelenumi{(\roman{enumi})}
\item
  \(\Longleftrightarrow\) (ii)：这由引理4 得出。
\item
  \(\Longrightarrow\) (iii)：显然。
\end{enumerate}

我们假设 k 是代数闭域，并证明 (iii) \(k\)\(\Longrightarrow\) (i)。先处理 \(\mathfrak{t} = \mathfrak{t}' = \mathfrak{g} = \mathfrak{sl}(2, k)\) 的情形。由于 \(\operatorname{ad}_{\mathfrak{g}} x\) 是幂零的，k² 的自同态 x 是幂零的（第 I 章第6节定理3），所以存在矩阵 \(x\)\(k^2\)\(A \in \mathbf{GL}(2, k)\)，使 \(AxA^{-1} = X\)，从而存在  的自同构 \(\alpha\)，使 \(\mathfrak{sl}(2, k)\)\(\alpha(x) = x'\)；现在 \(\alpha \in \operatorname{Aut}_e(\mathfrak{g})\)（§5，第3号，命题5的推论2）。现在转到一般情形；假设 \(\mathfrak{t}\) 和 \(\mathfrak{t}'\) G-共轭，并证明 x 和 \(x\)\(x'\) G-共轭。可以假设 \(\mathfrak{t} = \mathfrak{t}'\)。由上文，存在 \(\beta \in \operatorname{Aut}_e(\mathfrak{t})\)，使 \(\beta x = x'\)。现在，若 \(t \in \mathfrak{t}\) 满足 \(\operatorname{ad}_{\mathfrak{t}} t\) 幂零，则 \(\operatorname{ad}_{\mathfrak{g}} t\) 也幂零；所以 \(\beta\) 延拓为 \(\operatorname{Aut}_e(\mathfrak{g})\) 的一个元素。

注。只假设 \(k = k^2\) 时，命题1 的三个条件仍然等价（参见~习题1）。

\subsection{\texorpdfstring{\(\mathfrak{sl}_2\) 三元组在半单李代数中}{\textbackslash mathfrak\{sl\}\_2 三元组在半单李代数中}}

引理 5。令 V 为有限维向量空间，A、B 为 V 的自同态。假设 A 是幂零的且 \([A, [A, B]] = 0\)。则 AB 是幂零的。

置 C = {[}A, B{]}。由于 {[}A, C{]} = 0，

\[
[ \mathrm{A}, \mathrm{BC} ^ {p} ] = [ \mathrm{A}, \mathrm{B} ] \mathrm{C} ^ {p} = \mathrm{C} ^ {p + 1}
\]

对每个整数 \(p \geq 0\) 都成立。因此，\(\mathrm{Tr}(\mathrm{C}^{p}) = 0\) 对 \(p \geq 1\) 成立，这证明 C 是幂零的（《代数》第 VII 章第3节第5号命题13的推论4）。现在令 \(\bar{k}\) 为 k 的代数闭包，令 \(\lambda \in \bar{k}\)、\(x \in V \otimes_{k} \bar{k}\) 满足 \(\lambda x\)、\(x \neq 0\)。关系 \([[B, A], A] = 0\) 表明，对每个整数 \([B, A^{p}] = p[B, A]A^{p-1}\)，都有 \(p \geq 0\)。令 r 为满足 \(A^{r}x = 0\) 的最小整数。于是

\[
\lambda \mathrm{A} ^ {r - 1} x = \mathrm{A} ^ {r - 1} \mathrm{AB} x = \mathrm{A} ^ {r} \mathrm{B} x = \mathrm{BA} ^ {r} x - [ \mathrm{B}, \mathrm{A} ^ {r} ] x = - r [ \mathrm{B}, \mathrm{A} ] \mathrm{A} ^ {r - 1} x.
\]

由于 {[}B, A{]} 是幂零的且 \(A^{r-1}x \neq 0\)，可知 \(\lambda = 0\)。因此 AB 的所有特征值都为零，故引理得证。

引理 6。令 \(h, x \in \mathfrak{g}\) 满足 \([h, x] = 2x\)，且 \(h \in (\operatorname{ad} x)(\mathfrak{g})\)。则存在 \(y \in \mathfrak{g}\)，使得 \((x, h, y)\) 或为 \((0, 0, 0)\)，或为一个 \(\mathfrak{sl}_2\) 三元组。

令 \(\mathfrak{g}'\) 为可解李代数 \(kh + kx\)。由于 \(x \in [\mathfrak{g}', \mathfrak{g}']\)，\(\operatorname{ad}_{\mathfrak{g}} x\) 是幂零的（第 I 章第5节第3号定理1）；令 \(\mathfrak{n}\) 为其核。由于 {[}ad \(h\) , ad \(x] = 2\) ad \(x\)，有 (ad \(h\) ) \(\mathfrak{n} \subset \mathfrak{n}\)。令 \(z \in \mathfrak{g}\) 满足 \(h = -[x, z]\)。对任意整数 \(n \geq 0\)，置 \(\mathrm{M}_n = (\mathrm{ad}~x)^n\mathfrak{g}\)。若 \(n > 0\)，则有（§1，第1号，引理1）

\[
[ \mathrm{ad} z, (\mathrm{ad} x) ^ {n} ] = n ((\mathrm{ad} h) - n + 1) (\mathrm{ad} x) ^ {n - 1}
\]

所以，若 \(u \in M_{n-1}\)，

\[
n ((\operatorname{ad} h) - n + 1) u \in (\operatorname{ad} z) (\operatorname{ad} x) u + \mathrm{M} _ {n}.
\]

由于 (ad \(h)\mathfrak{n}\subset \mathfrak{n}\)，可得

\[
((\operatorname{ad} h) - n + 1) (\mathfrak {n} \cap \mathrm{M} _ {n - 1}) \subset \mathfrak {n} \cap \mathrm{M} _ {n}.
\]

由于 ad x 是幂零的，\(M_{n} = 0\) 对充分大的 n 成立。~因此，ad \(h|n\) 的特征值都是 \(\geq 0\) 的整数。于是 ad \(h + 2\) 在 n 上的限制可逆。

现在 \([h,z] + 2z\in \mathfrak{n}\)，因为

\[
\begin{array}{r} [ x, [ h, z ] + 2 z ] = [ [ x, h ], z ] + [ h, [ x, z ] ] + 2 [ x, z ] \\ = [ - 2 x, z ] + [ h, - h ] + 2 [ x, z ] = 0. \end{array}
\]

因此存在 \(z' \in \mathfrak{n}\)，使 \([h, z'] + 2z' = [h, z] + 2z\)，即置 \([h, y] = -2y\) 后有 \(y = z - z'\)。由于 \([x, y] = [x, z] = -h\)，证明完成。

命题 2（Jacobson–Morozov）。假设 \(\mathfrak{g}\) 是半单的。令 \(x\) 为 \(\mathfrak{g}\) 的非零幂零元素。存在 \(h, y \in \mathfrak{g}\)，使得 \((x, h, y)\) 是一个 \(\mathfrak{sl}_2\) 三元组。

令 \(\mathfrak{n}=\operatorname{Ker}(\operatorname{ad}x)^{2}\)。若 \(z\in\mathfrak{n}\)，则 {[}ad x, {[}ad x, ad z{]}{]} = ad({[}x, {[}x,z{]}{]}) = 0。由引理5，ad \(x\circ ad\) z 是幂零的，所以 \(\operatorname{Tr}(\operatorname{ad}x\circ\operatorname{ad}z)=0\)。这表明 x 关于 g 的 Killing 型 \(\Phi\) 与 n 正交。由于

\[
\Phi ((\operatorname{ad} x) ^ {2} y, y ^ {\prime}) = \Phi (y, (\operatorname{ad} x) ^ {2} y ^ {\prime})
\]

对所有 \(y, y' \in g\) 都成立，并且由于 \(\Phi\) 非退化，n 的正交补是 \((\operatorname{ad} x)^{2}\) 的像。因此存在 \(y' \in g\)，使 \(x = (\operatorname{ad} x)^{2}y'\)。置

\[
h = - 2 [ x, y ^ {\prime} ];
\]

我们有 \([h,x]=2x\) 且 \(h\in(\operatorname{ad}x)(\mathfrak{g})\)。现在只需应用引理6。

推论。假设 \(\mathfrak{g}\) 是半单的。令 \(\mathrm{G}\) 为包含 \(\mathfrak{g}\) 的 \(\operatorname{Aut}_e(\mathfrak{g})\) 的自同构群。将 g 中任意 \(\mathfrak{sl}_2\) 三元组 \((x,h,y)\) 与幂零元素 x 关联的映射在取商后，定义了从 \(sl_{2}\) 三元组的 G-共轭类到非零幂零元素的 G-共轭类的一个双射。

这由命题1和2 得出。

引理 7。令 K 为至少含有 4 个元素的交换域。令 G 为矩阵 \(\begin{pmatrix}\alpha&\beta\\0&\alpha^{-1}\end{pmatrix}\) 所成的群，其中 \(\alpha\in K^{*},\beta\in K\)。令 \(G'\) 为满足 \(\alpha=1\) 的这类矩阵所成的群。则 \(G'=(G,G)\)。

若 \(\alpha, \alpha' \in \mathrm{K}^*\) 且 \(\beta, \beta' \in \mathrm{K}\)，

\[
\begin{array}{c} \left( \begin{array}{c c} \alpha & \beta \\ 0 & \alpha^ {- 1} \end{array} \right) \left( \begin{array}{c c} \alpha^ {\prime} & \beta^ {\prime} \\ 0 & \alpha^ {- 1} \end{array} \right) \left( \begin{array}{c c} \alpha & \beta \\ 0 & \alpha^ {- 1} \end{array} \right) ^ {- 1} \left( \begin{array}{c c} \alpha^ {\prime} & \beta^ {\prime} \\ 0 & \alpha^ {- 1} \end{array} \right) ^ {- 1} \\ = \left( \begin{array}{c c} 1 & - \alpha^ {\prime} \beta^ {\prime} - \alpha \beta \alpha^ {2} + \alpha^ {2} \alpha^ {\prime} \beta^ {\prime} + \alpha \beta \\ 0 & 1 \end{array} \right). \end{array}
\]

特别地，

\[
\left( \begin{array}{c c} 1 & \beta \\ 0 & 1 \end{array} \right) \left( \begin{array}{c c} \alpha^ {\prime} & 0 \\ 0 & \alpha^ {- 1} \end{array} \right) \left( \begin{array}{c c} 1 & \beta \\ 0 & 1 \end{array} \right) ^ {- 1} \left( \begin{array}{c c} \alpha^ {\prime} & 0 \\ 0 & \alpha^ {- 1} \end{array} \right) ^ {- 1} = \left( \begin{array}{c c} 1 & \beta (1 - \alpha^ {2}) \\ 0 & 1 \end{array} \right).
\]

但存在 \(\alpha_0' \in \mathrm{K}^*\)，使 \(\alpha_0' \neq 1\) 且 \(\alpha_0' \neq -1\)，于是 \(k.(1 - \alpha_0'^2) = k\)，故引理得证。

命题 3。假设 \(\mathfrak{g}\) 是半单的。群 \(\mathrm{Aut}_e(\mathfrak{g})\) 等于其换位子群。若 \(\mathfrak{g}\) 可分裂，则 \(\mathrm{Aut}_e(\mathfrak{g})\) 是 \(\mathrm{Aut}_0(\mathfrak{g})\) 的换位子群。

令 \(x\) 为 \(\mathfrak{g}\) 的非零幂零元素。取 \(h, y \in \mathfrak{g}\)\((x, h, y)\)，使  为一个 \(\mathfrak{sl}_2\) 三元组（命题2）。由  生成的  的子代数 \(\mathfrak{s}\) 可识别为 。令  为  在 \(\mathfrak{g}\) 上的表示 \((x, h, y)\)\(\mathfrak{sl}(2, k)\)，并令  为与 \(\rho\) 相容的 \(z \mapsto \operatorname{ad}_{\mathfrak{g}} z\)\(\mathfrak{s} = \mathfrak{sl}(2, k)\)\(\mathfrak{g}\) 的表示（§1，第4号）。\(\pi\) 的像由 \(\mathbf{SL}(2, k)\)\(\rho\)\(\S 1\)\(\pi\) 的 \(\exp(t\operatorname{ad}_{\mathfrak{g}} x)\) 和 \(\exp(t\operatorname{ad}_{\mathfrak{g}} y)\) 生成（《代数》第 III 章第8节第9号命题17），因而包含于 \(t \in k\)\(\S 8\)\(\operatorname{Aut}_e(\mathfrak{g})\) 中。由于 \(\mathbf{SL}(2, k)\) 等于其换位子群（引理7及同上），\(\exp(\operatorname{ad}_{\mathfrak{g}} x)\) 属于 \(\operatorname{Aut}_e(\mathfrak{g})\) 的换位子群 G。因此 \(\operatorname{Aut}_e(\mathfrak{g})\) 等于 G。现在假设 \(\mathfrak{g}\) 可分裂。由于 \(\operatorname{Aut}_0(\mathfrak{g}) / \operatorname{Aut}_e(\mathfrak{g})\) 是交换的（§5，第3号，注3），上面的论证说明 \(\S 5\)\(\operatorname{Aut}_0(\mathfrak{g})\) 的换位子群是 \(\operatorname{Aut}_e(\mathfrak{g})\)。

\subsection{单元素}

定义 2。若存在 \(h\)，使  为 \(\mathfrak{g}\) 中的一个 \(x, y \in \mathfrak{g}\)\((x, h, y)\)\(\mathfrak{sl}_2\) 三元组，则称 \(\mathfrak{g}\) 的元素 h 为单元素。

我们也称 \(h\) 为 \(\mathfrak{sl}_2\) 三元组 \((x, h, y)\) 的单元素。

命题 4。令 \(h\) 为 \(\mathfrak{g}\) 的非零元素。则 \(h\) 是单元素，当且仅当存在 \(x \in \mathfrak{g}\)，使 \([h, x] = 2x\) 且 \(h \in (\operatorname{ad} x)(\mathfrak{g})\)。

必要性显然。充分性由引理6 得到。

命题 5。假设 g 是可分裂半单李代数。令 h 为 g 的分裂嘉当子代数，R 为 \((\mathfrak{g},\mathfrak{h})\) 的根的集合，B 为 R 的一个基。令 h 为属于 h 的 g 的一个单元素。则 h 在 \(\operatorname{Aut}_{e}(\mathfrak{g},\mathfrak{h})\) 下与 h 中满足对所有 \(h'\) 都有 \(\alpha(h') \in \{0,1,2\}\) 的元素 \(\alpha \in B\) 共轭。

\(\mathrm{ad}_{\mathfrak{g}}h\) 的特征值属于 \(\mathbf{Z}\)（§1，第2号，命题2的推论）。因此 \(h \in \mathfrak{h}_{\mathbf{Q}}\)。存在 \(w\) 的 Weyl 群中的元素 \((\mathfrak{g}, \mathfrak{h})\)，使对所有 \(\alpha(wh) \geq 0\) 都有 \(\alpha \in \mathrm{B}\)（第 VI 章第1节第5号定理2 (i)）。根据 §2 第2号定理2的推论，我们只需考虑对所有 \(\alpha(h) \in \mathbf{N}\) 都有 \(\alpha \in \mathrm{B}\) 的情形。令 \(\mathrm{R}_{+}\) 为关于 B 的正根集合，\(\mathrm{R}_{-} = -\mathrm{R}_{+}\)。\(\mathfrak{sl}_2\) 中存在形如 \(\mathfrak{g}\) 的 \((x, h, y)\) 三元组。令 T 为满足 \(\beta\) 的根 \(\beta(h) = -2\) 所成的集合。则 \(\mathrm{T} \subset \mathrm{R}_{-}\)，且 \(y \in \sum_{\beta \in \mathrm{T}} \mathfrak{g}^{\beta}\)。假设存在 \(\alpha \in \mathrm{B}\) 使 \(\alpha(h) > 2\)。对所有 \(\beta \in \mathrm{T}\)，有 \((\alpha + \beta)(h) > 0\)，所以 \(\alpha + \beta \notin \mathrm{R}_{-}\) 且 \(\alpha + \beta \neq 0\)；另一方面，由于 \(\beta \in \mathrm{R}_{-}\) 且 \(\alpha \in \mathrm{B}\)，有 \(\alpha + \beta \notin \mathrm{R}_{+}\)；因而 \(\alpha + \beta \notin \mathrm{R} \cup \{0\}\)，所以 \([\mathfrak{g}^{\alpha}, \mathfrak{g}^{\beta}] = 0\)。于是 \([y, \mathfrak{g}^{\alpha}] = 0\)。但 \(\mathrm{ad}_{\mathfrak{g}}y|\mathfrak{g}^{\alpha}\) 是单射，因为 \(\alpha(h) > 0\)（§1，第2号，命题2的推论）。这个矛盾证明了对所有 \(\alpha(h) \leq 2\) 都有 \(\alpha \in \mathrm{B}\)。

推论。若 k 是代数闭域且 \(k\)\(\mathfrak{g}\) 是秩为 \(l\) 的半单李代数，则相对于 ，\(\mathfrak{g}\) 的单元素共轭类至多有 \(\operatorname{Aut}_e(\mathfrak{g})\)\(3^l\) 个。

事实上，\(\mathfrak{g}\) 的每个半单元素都在 \(\mathrm{Aut}_e(\mathfrak{g})\) 下与 \(\mathfrak{h}\) 中的一个元素共轭。

引理 8。假设 k 是代数闭域且 \(k\)\(\mathfrak{g}\) 是半单的。令 \(h\) 为 \(\mathfrak{g}\) 的半单元素，且 \(\operatorname{ad} h\) 的特征值都是有理数。令 \(\mathfrak{g}^0 = \operatorname{Ker}(\operatorname{ad} h)\)，\(\mathfrak{g}^2 = \operatorname{Ker}(\operatorname{ad} h - 2)\)。令 \(G_h\) 为保持 h 不变的 \(\mathfrak{g}\) 的初等自同构所成的集合。令 \(h\)\(x \in \mathfrak{g}^2\) 满足 \([x, \mathfrak{g}^0] = \mathfrak{g}^2\)。则 \(G_h x\) 包含 \(\mathfrak{g}^2\) 的一个在 Zariski 拓扑下稠密且开的子集。

令 \(\mathfrak{h}\) 为 \(\mathfrak{g}^0\) 的一个嘉当子代数。这是包含 h 的 \(\mathfrak{g}\) 的一个嘉当子代数（第 VII 章第2节第3号命题10）。我们有 \(h\)\(h \in \mathfrak{h}_{\mathbf{Q}}\)。令 R 为 \((\mathfrak{g}, \mathfrak{h})\) 的根系，Q 为根权格。存在 R 的一个基 B，使对所有  都有 \(\alpha(h) \geq 0\)。\(\alpha \in B\)

令 \(\mathrm{U}\) 为满足对所有  都有  的 \(z \in \mathfrak{h}\) 所成的集合。令 \(\alpha(z) \neq 0\)\(\alpha \in \mathrm{B}\)\((H_{\alpha}')_{\alpha \in \mathrm{B}}\) 为与 B 对偶的 \(\mathfrak{h}\) 的基。若 \(\mathrm{B}\)\(z \in \mathrm{U}\)，则存在从 \(\mathrm{Q}\) 到 \(k^*\) 的同态，它将任意 \(\gamma \in \mathrm{Q}\) 映为 \(\prod_{\alpha \in \mathrm{B}} \alpha(z)^{\gamma(H_{\alpha}')}\)。根据 §5 命题2和4，在向量空间  上的自同态 \(\varphi(z)\) 在 \(\mathfrak{g}\)\(\mathfrak{g}^{\gamma}\) 上诱导比值为 \(\prod_{\alpha \in \mathrm{B}} \alpha(z)^{\gamma(H_{\alpha}')}\) 的伸缩，因此是 \(\mathfrak{g}\) 的初等自同构，并且显然属于 \(G_h\)。

令 \(s \in \mathfrak{h}\)。若 \(\gamma \in \mathrm{R}\) 满足 \(\mathfrak{g}^{\gamma} \cap \mathfrak{g}^{2} \neq 0\)，

\[
2 = \gamma (h) = \gamma \left(\sum_ {\alpha \in \mathrm{B}} \alpha (h) H _ {\alpha} ^ {\prime}\right) = \sum_ {\alpha \in \mathrm{B}} \alpha (h) \gamma (H _ {\alpha} ^ {\prime});
\]

由于对所有 \(\alpha(h) \geq 0\) 都有 \(\alpha \in \mathrm{B}\)，且 \(\gamma(H_{\alpha}^{\prime})\) 是全都 \(\geq 0\) 或全都 \(\leq 0\) 的整数，故对所有 \(\gamma(H_{\alpha}^{\prime}) \in \mathbf{N}\) 都有 \(\alpha \in \mathrm{B}\)。因此，我们可以（对 \(z \in \mathfrak{h}\)）考虑向量空间 \(\psi(z)\) 的自同态 \(\mathfrak{g}^2\)，它在 \(\mathfrak{g}^{\gamma} \cap \mathfrak{g}^2\) 上诱导比值为 \(\prod_{\alpha \in \mathrm{B}} \alpha(z)^{\gamma(H_{\alpha}^{\prime})}\) 的伸缩。映射 \(z \mapsto \psi(z)\) 从 \(\mathfrak{h}\) 到 \(\operatorname{End}(\mathfrak{g}^2)\) 是多项式的。对 \(z \in \mathrm{U}\)，有 \(\psi(z) = \varphi(z)|\mathfrak{g}^2\)。

令 \(\gamma_1, \ldots, \gamma_r\)\((\mathfrak{g}, \mathfrak{h})\)\(h\)\(y_1 \in \mathfrak{g}^{\gamma_1}, \ldots, y_r \in \mathfrak{g}^{\gamma_r}\)\(e^{\mathrm{ad} y_1} \ldots e^{\mathrm{ad} y_r} \in \mathrm{G}_h\)\(\rho\) 为在 h 上消失的  的互不相同的根。若 ，则 。于是可定义从 \(\mathfrak{h} \times \mathfrak{g}^{\gamma_1} \times \cdots \times \mathfrak{g}^{\gamma_r}\) 到 \(\mathfrak{g}^2\) 的映射 ，置

\[
\rho (z, y _ {1}, \dots , y _ {r}) = \psi (z) e ^ {\operatorname{ad} y _ {1}} \dots e ^ {\operatorname{ad} y _ {r}} x
\]

对 \(z \in \mathfrak{h}\)、\(y_1 \in \mathfrak{g}^{\gamma_1}, \ldots, y_r \in \mathfrak{g}^{\gamma_r}\)，此映射是多项式的，且 \(\rho(\mathrm{U}, \mathfrak{g}^{\gamma_1}, \ldots, \mathfrak{g}^{\gamma_r}) \subset \mathrm{G}_h x\)。根据第 VII 章附录 I 命题3和4，只需证明 \(\rho\) 的切线性映射在某一点满射。

现在令 T 为 \(z \mapsto \psi(z)\) 在 \(h_{0} = \sum_{\alpha \in B} H_{\alpha}'\) 处的切线性映射。于是 \(\mathrm{T}(z)\) 是 \(g^{2}\) 的自同态，它在 \(g^{\gamma} \cap g^{2}\) 上诱导比值为

\[
\sum_ {\alpha \in \mathrm{B}} \gamma (H _ {\alpha} ^ {\prime}) \alpha (h _ {0}) ^ {\gamma (H _ {\alpha} ^ {\prime}) - 1} \alpha (z) \prod_ {\beta \in \mathrm{B}, \beta \neq \alpha} \beta (h _ {0}) ^ {\gamma (H _ {\alpha} ^ {\prime})} = \sum_ {\alpha \in \mathrm{B}} \gamma (H _ {\alpha} ^ {\prime}) \alpha (z) = \gamma (z).
\]

因此，\(z \mapsto \rho(z,0,\ldots,0)\) 在 \(h_{0}\) 处的切线性映射是 \(z \mapsto [z,x]\)，其像为 \([x,h]\)。映射 \(y_{1} \mapsto \rho(h_{0},y_{1},0,\ldots,0)\) 在 0 处的切线性映射是 \(y_{1} \mapsto \psi(h_{0})[y_{1},x]\)；此映射的像为 \(\psi(h_{0})[x,\mathfrak{g}^{\gamma_{1}}] = [x,\mathfrak{g}^{\gamma_{1}}]\)。类似地，映射 \(y_{i} \mapsto \rho(h_{0},0,\ldots,0,y_{i},0,\ldots,0)\) 在 0 处的切线性映射的像为 \([x,\mathfrak{g}^{\gamma_{i}}]\)。最后，\(\rho\) 在 \((h_{0},0,\ldots,0)\) 处的切线性映射的像为

\[
[ x, \mathfrak {h} + \mathfrak {g} ^ {\gamma_ {1}} + \dots + \mathfrak {g} ^ {\gamma_ {r}} ] = [ x, \mathfrak {g} ^ {0} ] = \mathfrak {g} ^ {2}.\tag{Q.E.D.}
\]

*群 \(\mathrm{G}_h\) 是一个代数群，其李代数为 ad \(\mathfrak{g}^0_*\)

命题 6。假设 k 是代数闭域且 \(k\)\(\mathfrak{g}\)\(G\)\(\mathfrak{g}\) 是半单的。令 G 为包含 \(\operatorname{Aut}_e(\mathfrak{g})\) 的 \((x, h, y)\)\((x', h', y')\) 的自同构群。令  和  为  中的 \(\mathfrak{sl}_2\) 三元组。以下条件等价：\(\mathfrak{g}\)

\begin{enumerate}
\def\labelenumi{(\roman{enumi})}
\item
  h 与 \(h'\) G-共轭；
\item
  \((x,h,y)\) 与 \((x^{\prime},h^{\prime},y^{\prime})\) G-共轭。
\end{enumerate}

只需证明蕴含 (i) \(\Longrightarrow\) (ii)，且可立即归约到 \(h = h'\) 的情形。引入引理8中的 \(g^{2}\) 和 \(G_{h}\)。由 §1，第2号，命题2的推论，有 \(x \in g^{2}\) 且 \([x, g^{0}] = g^{2}\)。因此，\(G_{h}x\) 包含 \(g^{2}\) 的一个在 Zariski 拓扑下稠密且开的子集，\(G_{h}x'\) 也如此。所以存在 \(a \in G_{h}\)，使 \(a(x) = x'\)。我们有 \(a(h) = h\)，因而 \(a(y) = y'\)（第1号，引理1）。

推论 1。将任意 \(\mathfrak{sl}_2\) 三元组与其单元素关联的映射在取商后，定义了从 \(\mathfrak{sl}_2\) 三元组的 G-共轭类到单元素的 G-共轭类的一个双射。

推论 2。若 \(\operatorname{rk}(\mathfrak{g}) = l\)，则相对于 \(\operatorname{Aut}_e(\mathfrak{g})\)，\(\mathfrak{g}\) 的非零幂零元素的共轭类数至多为 \(3^l\)。

这由推论1、命题2的推论以及命题5的推论得到。

推论 3。若 \(\mathrm{rk}(\mathfrak{g}) = l\)，则相对于 \(\mathrm{Aut}_e(\mathfrak{g})\)，同构于 \(\mathfrak{g}\) 的 \(\mathfrak{sl}(2,k)\) 的子代数的共轭类数至多为 \(3^l\)。

这由推论1、命题1以及命题5的推论得到。

\subsection{主元素}

定义 3。假设 g 是半单的。

\begin{enumerate}
\def\labelenumi{(\roman{enumi})}
\item
  若 \(x\) 的幂零元素 \(\mathfrak{g}\) 的 \(\operatorname{Ker} \operatorname{ad} x\) 的维数等于 \(\mathfrak{g}\) 的秩，则称 x 为主幂零元素。
\item
  若 \(h\)\(\mathfrak{g}\) 的单元素 \(h\) 是正则的，且 \(\operatorname{ad} h\) 在 k 的代数闭包中的特征值属于 2Z，则称 h 为主单元素。\(k\)
\item
  若  的 \(\mathfrak{sl}_2\) 三元组 \((x,h,y)\) 在视为 \(\mathfrak{g}\)\(\mathfrak{g}\)\(kx + kh + ky\) 上的模时的长度等于 \(\mathfrak{g}\) 的秩，则称其为主三元组。
\end{enumerate}

命题 7。假设 g 是半单的。令 \((x,h,y)\) 为 g 中的一个 \(sl_{2}\) 三元组。以下条件等价：

\begin{enumerate}
\def\labelenumi{(\roman{enumi})}
\item
  \(x\) 是主幂零元素；
\item
  h 是主单元素；
\item
  \((x,h,y)\) 是主三元组。
\end{enumerate}

对 \(p \in Z\)，令 \(\mathfrak{g}^{p} = \operatorname{Ker}(\operatorname{ad} h - p)\)。令 \(g' = \sum_{p \in Z} g^{2p}\)。若将 g 视为 \(a = kx + kh + ky\) 上的模，则 \(g'\) 是奇数维单子模之和（§1，第2号，命题2的推论）。令 l（分别为 \(l'\)）为将 g（分别为 \(g'\)）视为 a-模时的长度。由 §1，第2号，

\[
\dim (\operatorname{Ker} \operatorname{ad} x) = l \geq l ^ {\prime} = \dim (\operatorname{Ker} \operatorname{ad} h) \geq \operatorname{rk} (\mathfrak {g}).
\]

(i) 与 (iii) 的等价性立即得到。另一方面，条件 (ii) 意味着 \(\dim(\operatorname{Ker} \operatorname{ad} h) = \operatorname{rk}(\mathfrak{g})\) 且 \(g' = g\)，换言之，

\[
\dim (\operatorname{Ker} \operatorname{ad} h) = \operatorname{rk} (\mathfrak {g})
\]

且 \(l = l'\)。条件 (ii) 与其他条件的等价性随即得到。

命题 8。假设 \(\neq 0\) 是非零半单李代数。令  为 g 的分裂嘉当子代数，R 为 \((\mathfrak{g},\mathfrak{h})\) 的根系，B 为 R 的一个基，令 \(h^{0}\) 为 h 中满足对所有  都有 \(\alpha(h^{0}) = 2\) 的元素。\(\alpha \in B\)

\begin{enumerate}
\def\labelenumi{(\roman{enumi})}
\item
  元素 \(h^0\) 是单元素且是主单元素。
\item
  \(x\)\(\mathfrak{g}\) 中满足存在形如  的 \(\mathfrak{sl}_2\) 三元组的元素 x，正是 \((x, h^0, y)\)\(\sum_{\alpha \in B} \mathfrak{g}^\alpha\) 中在每个 \(\mathfrak{g}^\alpha\) 中都有非零分量的元素。
\end{enumerate}

元素 \(h^0\) 就是 §7 第5号引理2 中考察的元素（参见同处公式 (1)）。由该引理可知，\(h^0\) 是主单元素，并且若 \(x \in \sum_{\alpha \in \mathrm{B}} \mathfrak{g}^\alpha\) 在每个 \(\mathfrak{g}^\alpha\) 中都有非零分量，则存在形如 \(\mathfrak{sl}_2\) 的 \((x, h^0, y)\) 三元组。反之，令 \((x, h^0, y)\) 为一个 \(\mathfrak{sl}_2\) 三元组。有 \([h^0, x] = 2x\)，所以 \(x \in \sum_{\gamma \in \mathrm{R}, \gamma(h^0) = 2} \mathfrak{g}^\gamma = \sum_{\alpha \in \mathrm{B}} \mathfrak{g}^\alpha\)。同样，\(y \in \sum_{\alpha \in \mathrm{B}} \mathfrak{g}^{-\alpha}\)。写

\[
h ^ {0} = \sum_ {\alpha \in \mathrm{B}} a _ {\alpha} H _ {\alpha} \quad \text { where } a _ {\alpha} > 0 \text { for all } \alpha \in \mathrm{B},
\]

\[
x = \sum_ {\alpha \in \mathrm{B}} X _ {\alpha} \qquad \text { where } X _ {\alpha} \in \mathfrak {g} ^ {\alpha} \text { for   all } \alpha \in \mathrm{B},
\]

\[
y = \sum_ {\alpha \in \mathrm{B}} X _ {- \alpha} \quad \text { where } X _ {- \alpha} \in \mathfrak {g} ^ {- \alpha} \text { for   all } \alpha \in \mathrm{B}.
\]

则

\[
\sum_ {\alpha \in \mathrm{B}} a _ {\alpha} H _ {\alpha} = h ^ {0} = [ y, x ] = \sum_ {\alpha , \beta \in \mathrm{B}} [ X _ {- \beta}, X _ {\alpha} ] = \sum_ {\alpha \in \mathrm{B}} [ X _ {- \alpha}, X _ {\alpha} ]
\]

所以对所有 \([X_{-\alpha}, X_{\alpha}] \neq 0\) 都有 \(\alpha \in \mathrm{B}\)。

推论。在可分裂半单李代数中，存在主幂零元素。

在非可分裂半单李代数中，0 可能是唯一的幂零元素。

命题 9。假设 \(k\) 是代数闭域且 \(\mathfrak{g}\) 是半单的。所有 \(\mathfrak{g}\) 的主单元素（分别为主幂零元素）都在 \(\operatorname{Aut}_e(\mathfrak{g})\) 下共轭。

沿用命题8的记号。令 h 为一个主单元素。它在 \(\operatorname{Aut}_{e}(\mathfrak{g})\) 下与 h 中的某个 \(h' \in h\) 共轭，使得对所有 \(\alpha(h') \in \{0,1,2\}\) 都有 \(\alpha \in B\)（第3号，命题5）。由于 \(h'\) 是主单元素，对所有 \(\alpha(h') \neq 0\) 都有 \(\alpha(h') \in 2Z\) 且 \(\alpha \in B\)，所以对所有 \(\alpha(h') = 2\) 都有 \(\alpha \in B\)，从而 \(h' = h^{0}\)。这证明了主单元素的断言。

令 \(x, x'\) 为主幂零元素。存在 \(\mathfrak{sl}_2\) 三元组 \((x, h, y)\)、\((x', h', y')\)。由命题7，\(h\) 和 \(h'\) 是主单元素，因而由上文它们在 \(\mathrm{Aut}_e(\mathfrak{g})\) 下共轭。所以 \(x\) 和 \(x'\) 在 \(\mathrm{Aut}_e(\mathfrak{g})\) 下共轭（命题6）。

引理 9。沿用命题8的记号。对 \(\mathfrak{g}^p = \mathrm{Ker}(\mathrm{ad}h^0 - p)\)，置 \(p \in \mathbf{Z}\)。令 \(\mathfrak{g}_*^2\) 为 \(\mathfrak{g}^2 = \sum_{\alpha \in \mathrm{B}} \mathfrak{g}^\alpha\) 中在每个 \(\mathfrak{g}^\alpha\) 中都有非零分量的元素所成的集合。令 \(\mathrm{R}_+\) 为关于 \(\mathrm{B}\) 的正根集合，\(\mathfrak{n}_+ = \sum_{\alpha \in \mathrm{R}_+} \mathfrak{g}^\alpha\)，并令 \(x \in \mathfrak{g}_*^2\)。则 \(e^{\mathrm{ad}\mathfrak{n}_+}. x = x + [\mathfrak{n}_+, \mathfrak{n}_+]\)。

显然 \(e^{\mathrm{ad} \mathfrak{n}_+}.x \subset x + [\mathfrak{n}_+, \mathfrak{n}_+]\)。我们证明，若 \(v \in [\mathfrak{n}_+, \mathfrak{n}_+]\)，则 \(x + v \in e^{\mathrm{ad} \mathfrak{n}_+}.x\)。置 \(\mathfrak{n}^{(p)} = \sum_{r \geq p} \mathfrak{g}^{2r}\)；只需证明

\[
x + v \in e ^ {\mathrm{ad} \mathfrak {n} _ {+}}. x + \mathfrak {n} ^ {(p)}
\]

对所有 \(p \geq 2\) 都成立。当 \(p = 2\) 时这显然成立，因为 \(\mathfrak{n}^{(2)} = [\mathfrak{n}_+, \mathfrak{n}_+]\)（§3，第3号，命题9 (iii)）。假设我们已经找到 \(z \in \mathfrak{n}_+\)，使 \(v + x - e^{\mathrm{ad}z}.x \in \mathfrak{n}^{(p)}\)。由于存在形如 \(\mathfrak{sl}_2\) 的 \((x, h^0, y)\) 三元组（命题8），§1 第2号命题2的推论说明 \([x, \mathfrak{g}^{2p-2}] = \mathfrak{g}^{2p}\)；因此存在 \(z' \in \mathfrak{g}^{2p-2} \subset \mathfrak{n}_+\)，使

\[
v + x - e ^ {\mathrm{ad} z}. x \in [ z ^ {\prime}, x ] + \mathfrak {n} ^ {(p + 1)}.
\]

所以 \(v + x \in e^{\mathrm{ad}(z + z')}.x + \mathfrak{n}^{(p + 1)}\)，归纳即证明了断言。

命题 10。假设 g 是半单的。令  为 g 的分裂嘉当子代数，R 为 \((\mathfrak{g},\mathfrak{h})\) 的根系，B 为 R 的一个基，\(R_{+}\) 为关于 B 的正根集合，且 \(n_{+}=\sum_{\alpha\in R_{+}}g^{\alpha}\)。属于 \(n_{+}\) 的主幂零元素正是 \(n_{+}\) 中在每个  对应的 \(g^{\alpha}\) 中都有非零分量的元素。\(\alpha\in B\)

命题8和引理9证明了这类元素是主幂零元素。我们证明逆命题。显然可以假设 \(\mathfrak{g}\) 是单的。令 \(h^0\) 和 \(\mathfrak{g}^p\) 如命题8和引理9中所述。令 \(\omega\) 为最高根，并置 \(\omega(h^0) = 2q\)；有 \(q = h - 1\)，其中 h 为 R 的 Coxeter 数，参见~第 VI 章第1节第11号命题31。于是 \(h\)\(\mathfrak{g}^{2q} = \mathfrak{g}^\omega\)、\(\mathfrak{g}^{-2q} = \mathfrak{g}^{-\omega}\)，且当  时 \(\mathfrak{g}^{2k} = 0\)。存在一个主 \(|k| > q\)\(\mathfrak{sl}_2\) 三元组 \((x^0, h^0, y^0)\)\((\mathrm{ad} x^0)^{2q}(\mathfrak{g}^{-\omega}) = \mathfrak{g}^\omega\)\((\mathrm{ad} x^0)^{2q} \neq 0\)\(x\)。由 §1 第2号命题2的推论，，所以 。令 x 为属于  的 \(\mathfrak{g}\) 的主幂零元素。若 \(\mathfrak{n}_+\)\(\bar{k}\) 是 k 的代数闭包，则 \(k\)\(x \otimes 1\) 和 \(x^0 \otimes 1\) 在 \(\mathfrak{g} \otimes_k \bar{k}\) 的某个自同构下共轭（命题9），所以 \((\mathrm{ad} x)^{2q} \neq 0\)。存在 \(\lambda \in \mathbb{R}\)，使 \((\mathrm{ad} x)^{2q} \mathfrak{g}^\lambda \neq 0\)。置 \(x = \sum_{n \geq 1} x_n\)，其中 \(x_n \in \mathfrak{g}^{2n}\)。则

\[
(\operatorname{ad} x) ^ {2 q} \mathfrak {g} ^ {\lambda} \subset (\operatorname{ad} x _ {1}) ^ {2 q} \mathfrak {g} ^ {\lambda} + \sum_ {k > 4 q + \lambda (h ^ {0})} \mathfrak {g} ^ {k} = (\operatorname{ad} x _ {1}) ^ {2 q} \mathfrak {g} ^ {\lambda},
\]

由于 \(4q + \lambda(h^0) \geq 4q - 2q = 2q\)。现在，当 \((\operatorname{ad} x_1)^{2q} \mathfrak{g}^\lambda \subset \mathfrak{g}^{4q + \lambda(h^0)}\) 时，\(\lambda = -\omega\)。因此，\((\operatorname{ad} x_1)^{2q} \mathfrak{g}^{-\omega} = \mathfrak{g}^\omega\)。我们有 \(\omega = \sum_{\alpha \in \mathrm{B}} n_\alpha \alpha\)，其中对所有 \(n_\alpha > 0\) 都有 \(\alpha \in \mathrm{B}\)（第 VI 章第1节第8号注）。若存在 \(\alpha_0 \in \mathrm{B}\)，使 \(x_1 \in \sum_{\alpha \in \mathrm{B}, \alpha \neq \alpha_0} \mathfrak{g}^\alpha\)，则关系

\[
\omega \notin - \omega + \sum_ {\alpha \in \mathrm{B}, \alpha \neq \alpha_ {0}} k \alpha
\]

这意味着对所有 p 都有 \(\mathfrak{g}^{\omega} \not\subset (\operatorname{ad} x_1)^p \mathfrak{g}^{-\omega}\)\(p\)\(x_1\)，这是荒谬的。因此，x₁ 在每个 \(\mathfrak{g}^{\alpha}\) 中的分量对所有 \(\alpha \in B\) 都非零。

\section{Chevalley 序}

\subsection{格与序}

令 V 为一个 Q-向量空间。V 中的格是 V 的一个自由 Z-子模 V，使得由 V 嵌入 V 所诱导的 Q-线性映射 \(\alpha_{\mathcal{V},V} : \mathcal{V} \otimes_{Z} Q \to V\) 为双射。当 V 有限维时，这等价于说 V 是有限型的 Z-子模并生成 Q-向量空间 V（回忆一下，有限型的无挠 Z-模是自由的，见《代数》第 VII 章第4节第4号定理的推论2）；此外，此时我们的定义是《交换代数》第 VII 章第4节第1号定义1的特例（同上，例3）。若 W 是 V 的向量子空间，且 V 是 V 中的一个格，则 \(V \cap W\) 是 W 中的格。

若 V 是 Q-代数，则 V 中的一个序是底层向量空间中的格 V，且是 V 的 Z-子代数；此时映射 \(\alpha_{V,V}\) 是 Q-代数的同构。若 V 是有单位的 Q-代数，则 V 中的幺序是包含单位元的 V 中的序。

假设 V 是一个 Q-双代数，余积为 c，余单位为 \(\gamma\)。若 V 是向量空间 V 中的一个格，则典则映射 \(i: V \otimes_{Z} V \to V \otimes_{Q} V\) 是单射；V 中的双序是有单位代数 V 中的一个幺序 V，使得 \(\gamma(\mathcal{V}) \subset \mathbf{Z}\) 且 \(c(\mathcal{V}) \subset i(\mathcal{V} \otimes_{\mathbf{Z}} \mathcal{V})\)；映射

\[
\gamma_ {\mathcal {V}} \colon \mathcal {V} \to \mathbf {Z} \text { and } c _ {\mathcal {V}} \colon \mathcal {V} \to \mathcal {V} \otimes_ {\mathbf {Z}} \mathcal {V}
\]

由 \(\gamma\) 和 c 诱导的  及  赋予 V 以 Z-双代数结构，而映射 \(\alpha_{\mathcal{V},V}\) 此时是 Q-双代数的同构。

\subsection{双代数中的除幂}

令 A 为有单位 k-代数，\(x \in A\)，\(d \in k\)，\(n \in N\)。置

\[
x ^ {(n, d)} = \frac {x (x - d) \dots (x - d (n - 1))}{n !} = \prod_ {i = 0} ^ {n - 1} (x - i d) / (i + 1).\tag{1}
\]

特别地，\(x^{(0,d)} = 1\)，\(x^{(1,d)} = x\)。约定当 n 为小于 0 的整数时 \(x^{(n,d)} = 0\)。置\(n\)\(< 0\)

\[
x ^ {(n)} = x ^ {(n, 0)} = \frac {x ^ {n}}{n !}\tag{2}
\]

\[
\binom {x} {n} = x ^ {(n, 1)} = \frac {x (x - 1) \ldots (x - n + 1)}{n !}.\tag{3}
\]

命题 1。令 A 为双代数，余积为 \(c\)，\(x\) 为 A 的一个本原元（第 II 章第1节第2号）。则

\[
c (x ^ {(n, d)}) = \sum_ {p \in \mathbf {N}, q \in \mathbf {N}, p + q = n} x ^ {(p, d)} \otimes x ^ {(q, d)}.\tag{4}
\]

当 \(n \leq 0\) 时命题显然成立。我们对 \(n\) 作归纳。若公式 (4) 对 \(n\) 成立，则

\[
\begin{array}{l} (n + 1) c (x ^ {(n + 1, d)}) = c (x - d n) c (x ^ {(n, d)}) \\ \quad = (x \otimes 1 + 1 \otimes x - d n   1 \otimes 1) c (x ^ {(n, d)}) \\ \quad = \sum_ {p + q = n} [ x x ^ {(p, d)} \otimes x ^ {(q, d)} + x ^ {(p, d)} \otimes x x ^ {(q, d)} - (p + q) d x ^ {(p, d)} \otimes x ^ {(q, d)} ] \\ \quad = \sum_ {p + q = n} (x - p d) x ^ {(p, d)} \otimes x ^ {(q, d)} + \sum_ {p + q = n} x ^ {(p, d)} \otimes (x - q d) x ^ {(q, d)} \\ \quad = \sum_ {p + q = n} (p + 1) x ^ {(p + 1, d)} \otimes x ^ {(q, d)} + \sum_ {p + q = n} (q + 1) x ^ {(p, d)} \otimes x ^ {(q + 1, d)} \\ \quad = \sum_ {r + s = n + 1} r x ^ {(r, d)} \otimes x ^ {(s, d)} + \sum_ {r + s = n + 1} s x ^ {(r, d)} \otimes x ^ {(s, d)} \\ \quad = (n + 1) \sum_ {r + s = n + 1} x ^ {(r, d)} \otimes x ^ {(s, d)}, \end{array}
\]

因而公式 (4) 对 \(n + 1\) 成立。

\subsection{POINCAR\'E-BIRKHOFF-WITT 定理的整形式}

令 g 为有限维 Q-李代数，\(U(\mathfrak{g})\) 为其包络双代数。若 I 为全序集，\(\mathbf{x} = (x_i)_{i \in \mathrm{I}}\) 为 g 中元素的一个族，\(\mathbf{n} = (n_i)_{i \in \mathrm{I}} \in \mathbf{N}^{(\mathrm{I})}\) 为一个多重指标，置

\[
\mathbf {x} ^ {(\mathbf {n})} = \prod_ {i \in \mathrm{I}} \frac {x _ {i} ^ {n _ {i}}}{n _ {i} !},\tag{5}
\]

其中乘积在 \(\mathrm{U}(\mathfrak{g})\) 中按照全序集 I 计算。

定理 1。令 \(\mathcal{U}\) 为双代数 \(\mathrm{U}(\mathfrak{g})\) 中的双序。令 \(\mathcal{G} = \mathcal{U} \cap \mathfrak{g}\)，这是李代数 \(\mathfrak{g}\) 中的一个序。令 \((x_i)_{i \in \mathrm{I}}\) 为 \(\mathcal{G}\) 的一个基。给 I 赋予一个全序，并假设对每个 \(\mathbf{n} \in \mathbf{N}^{\mathrm{I}}\)，给定 { 中的一个元素 }[{n}]\(\mathcal{U}\)，使 {[}n{]} - \(\mathbf{x}^{(\mathbf{n})}\) 在 \(< |\mathbf{n}|\) 中的滤子次数小于 \(\mathrm{U}(\mathfrak{g})\)。则 { 的 }[{n}]\(\mathbf{n} \in \mathbf{N}^{\mathrm{I}}\) 所成的族是 Z-模 \(\mathcal{U}\) 的一个基。

对 \(p \in N\)，令 \(U_{p}(\mathfrak{g})\) 为 \(U(\mathfrak{g})\) 中滤子次数 \(\leq p\) 的元素所成的集合；则 \(U_{p}(\mathfrak{g}) / U_{p-1}(\mathfrak{g})\) 的 \(\mathbf{x}^{(\mathbf{n})}\) 在 \(|n| = p\) 中的像构成此 Q-向量空间的一个基（第 I 章第2节第7号定理1）；因此 {[}n{]} 构成 Q-向量空间 \(U(\mathfrak{g})\) 的一个基。还需证明以下断言（其中置 \(M = N^{I}\)）：

\((^{*})\) 若 \(u\in \mathcal{U}\)、\((a_{\mathbf{n}})\in \mathbf{Z}^{(\mathrm{M})}\) 且 \(d\in \mathbf{N} - \{0\}\) 满足

\[
d u = \sum_ {\mathbf {n} \in \mathrm{M}} a _ {\mathbf {n}} [ \mathbf {n} ],\tag{6}
\]

则 \(d\) 整除每个 \(a_{\mathbf{n}}\)。

对每个整数 \(r \geq 0\)，引入迭代余积

\[
c _ {i}: \mathcal {U} \rightarrow \mathbf {T} ^ {r} (\mathcal {U}) = \mathcal {U} \otimes \mathcal {U} \otimes \dots \otimes \mathcal {U};
\]

按定义，\(c_0\) 是 \(\mathcal{U}\) 的余单位，\(c_1 = \mathrm{Id}_{\mathcal{U}}\)，\(c_2 = c\)（\(\mathcal{U}\) 的余积），并且当 \(r \geq 2\) 时，\(c_{r+1}\) 定义为复合映射 \(p \circ (c_r \otimes 1) \circ c\)：

\[
\mathcal {U} \stackrel {{c}} {{\longrightarrow}} \mathcal {U} \otimes_ {\mathbf {Z}} \mathcal {U} \stackrel {{c _ {r} \otimes 1}} {{\longrightarrow}} \mathbf {T} ^ {r} (\mathcal {U}) \otimes_ {\mathbf {Z}} \mathcal {U} \stackrel {{p}} {{\longrightarrow}} \mathbf {T} ^ {r + 1} (\mathcal {U})
\]

其中 p 由利用代数 \(\mathbf{T}(\mathcal{U})\) 中的乘法来定义。此外，考虑 U 到 \(\pi\) 的典则投影 \(U^{+} = Ker c_{0}\)，以及复合映射

\[
c _ {r} ^ {+} = \mathbf {T} ^ {r} (\pi) \circ c _ {r}: \mathcal {U} \to \mathbf {T} ^ {r} (\mathcal {U} ^ {+}).
\]

引理 1。令 \(\mathbf{n} \in \mathbf{N}^{\mathrm{I}}\)。若 \(|\mathbf{n}| < r\)，则 \(c_{r}^{+}([ \mathbf{n}]) = 0\)。若 \(|\mathbf{n}| = r\)，则

\[
c _ {r} ^ {+} ([ \mathbf {n} ]) = \sum_ {\varphi} x _ {\varphi (1)} \otimes x _ {\varphi (2)} \otimes \dots \otimes x _ {\varphi (r)},\tag{7}
\]

其中 \(\varphi\) 属于从 \(\{1,2,\ldots,r\}\) 到 I 的映射所成的集合，并且对每个值 \(i\in I\)，它取值 i 恰好 \(n_{i}\) 次。

由命题1，

\[
c _ {r} \left(\mathbf {x} ^ {(\mathbf {n})}\right) = \sum \mathbf {x} ^ {\left(\mathbf {p} _ {1}\right)} \otimes \dots \otimes \mathbf {x} ^ {\left(\mathbf {p} _ {r}\right)}
\]

其中求和遍历由 M 中 r 个元素组成的序列 \((\mathbf{p}_{1},\ldots,\mathbf{p}_{r})\) 所成的集合，这些元素满足 \(p_{1}+\cdots+p_{r}=n\)。根据第 II 章第1节第3号命题6，将 \(c_{r}^{+}\) 通过线性延拓为从 \(\mathrm{U}(\mathfrak{g})\) 到 \(\mathbf{T}^{r}(\mathrm{U}^{+}(\mathfrak{g}))\) 的映射后，它在 \(\mathrm{U}_{r-1}(\mathfrak{g})\) 上为零。因此，对 \(r\geq|\mathbf{n}|\)，

\[
c _ {r} ^ {+} ([ \mathbf {n} ]) = c _ {r} ^ {+} (\mathbf {x} ^ {(\mathbf {n})}) = \sum \pi (\mathbf {x} ^ {(\mathbf {p} _ {1})}) \otimes \dots \otimes \pi (\mathbf {x} ^ {(\mathbf {p} _ {r})}).\tag{8}
\]

当 \(r > |n|\) 时，关系 \(p_{1} + \cdots + p_{r} = n\) 蕴含至少有一个 \(p_{i}\) 为零，所以 \(c_{r}^{+}([n]) = 0\)。当 \(r = |n|\) 时，(8) 的第三项中非零的项只有满足 \(|p_{1}| = \cdots = |p_{r}| = 1\) 的项，因而得到 (7)。

回到定理1的证明。沿用 (*) 的记号，并对 \(|n|\) 作降阶归纳，证明 d 整除 \(a_{n}\)，当 \(|n|\) 充分大时这显然成立。若对 \(a_{n}\) 有 d 整除 \(|n| > r\)，则置

\[
u ^ {\prime} = u - \sum_ {| \mathbf {n} | > r} (a _ {\mathbf {n}} / d) [ \mathbf {n} ] \in \mathcal {U},
\]

有

\[
d u ^ {\prime} = \sum_ {| \mathbf {n} | \leq r} a _ {\mathbf {n}} [ \mathbf {n} ].\tag{9}
\]

对从 \(\varphi\) 到 I 的任意映射 \(\{1,2,\ldots,r\}\)，置

\[
e _ {\varphi} = x _ {\varphi (1)} \otimes \dots \otimes x _ {\varphi (r)}
\]

并令 \(a_{\varphi} = a_{\mathbf{n}}\)，其中 \(\mathbf{n} = (\operatorname{Card}\varphi^{-1}(i))_{i\in \mathrm{I}}\)。由引理1，(9) 蕴含

\[
d c _ {r} ^ {+} (u ^ {\prime}) = \sum_ {\varphi \in \mathrm{I} ^ {r}} a _ {\varphi} e _ {\varphi}\tag{10}
\]

所以 \(c_r^+(u') \in \mathbf{T}^r(\mathcal{U}^+) \cap \mathbf{Q}\mathbf{T}^r(\mathcal{G})\)。但是 \(\mathcal{G}\) 作为 \(\mathcal{U}^+\) 的子模是一个直因子（《代数》第 VII 章第4节第3号定理1的推论），所以子模 \(\mathbf{T}^r(\mathcal{G})\) 是 \(\mathbf{T}^r(\mathcal{U}^+)\) 的直因子，从而 \(c_r^+(u') \in \mathbf{T}^r(\mathcal{G})\)。另一方面，由假设，\(x_i\) 构成 \(\mathcal{G}\) 的一个基，所以 \(e_\varphi\) 构成 \(\mathbf{T}^r(\mathcal{G})\) 的一个基。于是 (10) 证明 d 整除 \(d\)\(a_\varphi\)，即当  时 d 整除 \(a_{\mathbf{n}}\)。这证明了 (*)。\(|\mathbf{n}| = r\)